% Layout reconstructed from OpenAI's Finite time blowup for Navier--Stokes
% (8 September 2026); the public release does not include a TeX template.
% Reference PDF: https://cdn.openai.com/pdf/32d9f210-8b73-45e0-91bc-82a30aef8a9a/navier-stokes.pdf
\documentclass[11pt,letterpaper,reqno]{amsart}
\usepackage[margin=30mm]{geometry}
\usepackage[T1]{fontenc}
\usepackage{amsmath,amssymb,mathtools}
\numberwithin{equation}{section}
\usepackage[sc]{mathpazo}
\usepackage{microtype}
\usepackage{etoolbox}
% Match the compact first-page title placement in the reference PDF.
\makeatletter
\patchcmd{\@maketitle}{\global\topskip42\p@\relax}{\global\topskip0\p@\relax}{}{
  \PackageError{DR-Yaglom}{Cannot adjust title placement}{Check the amsart version.}}
\makeatother
\usepackage{enumitem}
\usepackage[colorlinks=true,linkcolor=blue,citecolor=blue,urlcolor=blue]{hyperref}
\hypersetup{
  pdftitle={Exact critical asymptotics for Derrida--Retaux recursions},
  pdfauthor={Ruiqi Ding; Zehua He; Yutao Liang; Dong Wang; Yushu Zheng}
}
\theoremstyle{plain}
\newtheorem{theorem}{Theorem}[section]
\newtheorem{proposition}[theorem]{Proposition}
\newtheorem{lemma}[theorem]{Lemma}
\newtheorem{corollary}[theorem]{Corollary}
\theoremstyle{remark}
\newtheorem*{remark}{Remark}
\newcommand{\E}{\mathbb E}\newcommand{\Pp}{\mathbb P}
\newcommand{\Zp}{\mathbb Z_{\ge0}}
\newcommand{\dd}{\,\mathrm d}
\newcommand{\norm}[1]{\left\lVert#1\right\rVert}
\newcommand{\floor}[1]{\lfloor#1\rfloor}
\DeclareMathOperator{\supp}{supp}

\title[Critical Derrida--Retaux recursions]{Exact critical asymptotics for Derrida--Retaux recursions}
\DeclareRobustCommand{\authormark}[1]{\textsuperscript{\MakeLowercase{#1}}}
\author[R. Ding]{Ruiqi Ding\authormark{1,2,a}}
\author[Z. He]{Zehua He\authormark{1,2,b}}
\author[Y. Liang]{Yutao Liang\authormark{1,2,c}}
% Dong Wang: https://arxiv.org/html/2512.20343v1
\author[D. Wang]{Dong Wang\authormark{2,d}}
\author[Y. Zheng]{Yushu Zheng\authormark{3,e}}
\newcommand{\authoremail}[2]{\textsuperscript{#1}\href{mailto:#2}{\textnormal{#2}}}
\newcommand{\frontaffiliations}{%
  \par\vspace{8pt}
  \begingroup
  \normalfont\footnotesize\centering
  \hypersetup{urlcolor=black}
  \textsuperscript{1}\textit{State Key Laboratory of Mathematical Sciences,
  Academy of Mathematics and Systems Science, Chinese Academy of Sciences}\par
  \authoremail{a}{dingruiqi@amss.ac.cn},
  \authoremail{b}{hezehua@amss.ac.cn},
  \authoremail{c}{liangyutao@amss.ac.cn}\par\vspace{5pt}
  \textsuperscript{2}\textit{School of Mathematical Sciences,
  University of Chinese Academy of Sciences}\par
  \authoremail{d}{wangdong@wangd-math.xyz}\par\vspace{5pt}
  \textsuperscript{3}\textit{School of Mathematics,
  Shanghai University of Finance and Economics}\par
  \authoremail{e}{yszheng666@gmail.com}\par
  \endgroup}
\makeatletter
\apptocmd{\@setauthors}{\frontaffiliations}{}{
  \PackageError{DR-Yaglom}{Cannot place affiliations on the first page}{Check the amsart version.}}
\makeatother
\date{}
\begin{document}
\begin{abstract}
We study the critical Derrida--Retaux recursion with a fixed number $m\ge2$ of independent inputs and a nonnegative integer-valued initial law.
%
Assuming $\E[X_0^3m^{X_0}]<\infty$ and excluding the deterministic binary fixed point, we prove that the survival probability is asymptotic to $4/((m-1)^2n^2)$ and identify the locally uniform limit of $n^2\rho_n(\floor{nx})$, including at $x=0$, where $\rho_n(j)=(m-1)m^{j+1}\Pp(X_n=j+1)/\E[m^{X_n}]$.
%
We also obtain exact asymptotics for $\E[m^{X_n}]-1$, its one-step decrement and $\Pp(X_n=k)$ for each fixed $k\ge1$, together with convergence of the conditional law to a geometric distribution with total variation error $O(\log n/n)$.
%
The proof combines a cubic moment identity, a truncated convolution estimate and Kotani's uniqueness theorem for strings of entrance type.
%
A quantitative third-moment tail bound determines the small-time asymptotics needed to apply this theorem without a fourth-moment assumption.
%
For each fixed $r\in[0,3)$, the assumption $\E[X_0^rm^{X_0}]<\infty$ does not, in general, guarantee this survival asymptotic.
\end{abstract}
\maketitle

\section{Introduction}
We consider the Derrida--Retaux model on an $m$-ary tree, where $m\ge2$ is a fixed integer. For each $n\ge0$, we take such a tree with all its leaves at distance $n$ from the root. The leaves carry independent nonnegative integer-valued masses, all with the same law as a given random variable $X_0$.

We assign masses to the remaining vertices by working from the leaves towards the root. Once the masses at the $m$ children of a vertex are known, we add them, subtract one unit, and assign the result to that vertex if it is positive; otherwise, we assign zero. Repeating this operation gives a mass at the root, which we denote by $X_n$. Thus $n$ counts the number of steps from the leaves to the root. In a tree of height $n+1$, the masses at the children of the root are independent copies of $X_n$: each is obtained from a tree of height $n$, and these trees depend on disjoint sets of independent initial masses. This gives the recurrence relation \eqref{eq:model} below.

The binary model was studied by Collet, Eckmann, Glaser and Martin~\cite{CEGM} in connection with a spin glass model, and Derrida and Retaux~\cite{DR} later used it as a simplified hierarchical model for the depinning transition. The behavior of the mass at the root depends on the initial law. Under $\E[X_0m^{X_0}]<\infty$, its mean can grow exponentially with the height of the tree or remain bounded. For integer-valued initial masses, the critical case separating these regimes is characterized by
\[
 (m-1)\E[X_0m^{X_0}]=\E[m^{X_0}].
\]
This criterion is due to Collet et al.~\cite{CEGM}; see also \cite[Theorem~1]{CDHLS} for general $m$.

This paper concerns the critical case. We first study the survival probability $\Pp(X_n>0)$, namely, the probability that some positive mass remains at the root after $n$ steps of the construction. Our aim is to determine how this probability decays as $n\to\infty$, including its leading constant. We also study the conditional probabilities $\Pp(X_n=k\mid X_n>0)$, $k\ge1$. We seek the limiting conditional distribution as the height of the tree tends to infinity. This is often called a \emph{Yaglom-type theorem}, after Yaglom's classical result~\cite{Yaglom} on the distribution of a branching process conditioned on non-extinction. In the classical critical Galton--Watson setting, the limit is exponential after rescaling the population size; in the present model, we obtain a geometric limit for $X_n$ without rescaling.

For nonconstant critical initial laws with suitable moment assumptions, the predictions formulated in \cite[Conjectures~1 and~2]{CDHLS} are
\[
 \Pp(X_n>0)\sim\frac4{(m-1)^2n^2},
 \qquad
 \Pp(X_n=k\mid X_n>0)\longrightarrow(m-1)m^{-k}
 \quad(k\ge1).
\]
The first prediction specifies how often positive mass remains at the root of a large tree, including the leading constant. The second says that, conditional on this event, the remaining mass approaches a geometric distribution on the positive integers.

Several rigorous estimates support these predictions. Chen, Derrida, Hu, Lifshits and Shi~\cite{CDHLS} proved that $\E[m^{X_n}]-1$ is of order $n^{-1}$ when $\E[X_0^3m^{X_0}]<\infty$. Under the stronger assumption $\E[s^{X_0}]<\infty$ for some $s>m$, Chen, Hu and Shi~\cite[Theorem~1.2]{CHS} proved that both $\Pp(X_n>0)$ and $\E[X_n]$ are $n^{-2+o(1)}$. Their proof studies pivotal vertices and open paths in the tree and uses a coupling argument. In particular, it establishes the predicted exponent two for the survival probability. An estimate of the form $n^{-2+o(1)}$, however, does not determine a leading constant, so further estimates are needed for the precise asymptotic above.

There are also related models that are exactly solvable for certain families of initial distributions~\cite{HMP,LZ,AHM}. The continuous coalescence equation studied by Derrida and Retaux~\cite[Section~3.2]{DR} and Chen, Dagard, Derrida and Shi~\cite[Section~2]{CDDS} has explicit exponential solutions at criticality, which suggest a possible scaling limit for the discrete model. Derrida and Shi~\cite{DS} review these predictions and explain the role of moment assumptions on the initial law. To use the suggested limiting function to compute probabilities for the discrete model, one must prove convergence with enough control near the small positive values of $X_n$.

In this paper, we prove the survival asymptotic and the conditional geometric limit above for every fixed $m\ge2$ and every critical initial law satisfying $\E[X_0^3m^{X_0}]<\infty$, except for the deterministic law $X_0=1$ when $m=2$. In the exceptional case the mass remains identically $1$ at every generation. We also obtain exact asymptotics for $\Pp(X_n=k)$ for each fixed $k\ge1$, for $\E[X_n]$, and for $\E[m^{X_n}]-1$, together with a total variation error of order $O(\log n/n)$ for the conditional geometric approximation. We also show that, for any fixed $r\in[0,3)$, the assumption $\E[X_0^rm^{X_0}]<\infty$ does not guarantee the survival asymptotic.

The main step is to prove that $n^2\rho_n(\floor{nx})$ converges to $4e^{-2x}$ uniformly on every bounded interval $[0,R]$, where $\rho_n$ is defined below from the positive part of the exponentially tilted law of $X_n$. Since $\Pp(X_n=k)$ is expressed in terms of $\rho_n(k-1)$, each fixed $k\ge1$ corresponds to $x=(k-1)/n\to0$. The boundary limit $n^2\rho_n(0)\to4$, together with bounds on $|\rho_n(j+1)-\rho_n(j)|$ and $\E[m^{X_n}]\to1$, then gives the exact asymptotics of $\Pp(X_n=k)$ and $\Pp(X_n>0)$. We give the precise definitions and statements next, followed by an outline of the proof.

\subsection{Model and main result}
Fix an integer $m\ge2$. Starting from a nonnegative integer-valued random variable $X_0$, define the laws of $X_n$ recursively by
\begin{equation}\label{eq:model}
 X_{n+1}\overset{\mathrm d}=\bigl(X_n^{(1)}+\cdots+X_n^{(m)}-1\bigr)_+,
\end{equation}
where $x_+=\max\{x,0\}$ and, at each step, the variables $X_n^{(1)},\ldots,X_n^{(m)}$ are independent copies of $X_n$. We assume
\begin{equation}\label{eq:assumptions}
 (m-1)\E[X_0m^{X_0}]=\E[m^{X_0}]<\infty,
 \qquad \E[X_0^3m^{X_0}]<\infty.
\end{equation}
The first equality is the criticality condition of Collet et al.~\cite{CEGM}; see also \cite[Theorem~1]{CDHLS} for general $m$. The second is the moment condition used throughout the proof.
When $m=2$, we exclude $\Pp(X_0=1)=1$. This is a deterministic fixed point of \eqref{eq:model}. No other critical deterministic law exists.

Write
\begin{equation}\label{eq:notation}
 p_n:=\Pp(X_n>0),\quad G_n:=\E[m^{X_n}],\quad \varepsilon_n:=G_n-1,\quad 
 q_{n,k}:=\frac{m^k\Pp(X_n=k)}{G_n},
\end{equation}
and define
\begin{equation}\label{eq:rho}
 \rho_n(j):=(m-1)q_{n,j+1},\qquad j\in\Zp.
\end{equation}
The law $q_n=(q_{n,k})_{k\ge0}$ is obtained by tilting the law of $X_n$: we multiply $\Pp(X_n=k)$ by $m^k$ and divide by $G_n$. For each $j\ge0$, $\rho_n(j)$ is $m-1$ times the probability assigned to $j+1$ by the tilted law $q_n$. In particular,
\[
 \rho_n(0)=\frac{m(m-1)}{G_n}\Pp(X_n=1).
\]
Thus, $j=0$ corresponds to $X_n = 1$ in the original probability.

Criticality is preserved by \eqref{eq:model}, as shown in \eqref{eq:qmean}, and the factor $m-1$ gives the normalization
\[
 \sum_{k\ge0}q_{n,k}=1,\qquad
 \sum_{j\ge0}(j+1)\rho_n(j)=1.
\]
The sequence $\rho_n$ itself need not have total mass one. We use it to study $q_{n,k}$ at positive indices proportional to $n$. Its values near $j=0$ will then determine the leading terms of the original probabilities at fixed positive integers.

\begin{theorem}\label{thm:profile}
Assume \eqref{eq:assumptions} and exclude the deterministic binary law specified above. For every $R\in[0,\infty)$,
\begin{equation}\label{eq:mainprofile}
 \sup_{0\le x\le R}\left|n^2\rho_n(\floor{nx})-4e^{-2x}\right|\longrightarrow0.
\end{equation}
Consequently,
\begin{align}
 n(G_n-1)&\longrightarrow\frac2{m-1}, &
 n^2p_n&\longrightarrow\frac4{(m-1)^2},\label{eq:main1}\\
 n^2(\varepsilon_n-\varepsilon_{n+1})&\longrightarrow\frac2{m-1}, &
 n^2\E[X_n]&\longrightarrow\frac{4m}{(m-1)^3}.\label{eq:main2}
\end{align}
For every fixed $k\ge1$,
\begin{equation}\label{eq:atoms}
 n^2\Pp(X_n=k)\longrightarrow \frac4{(m-1)m^k},\qquad
 \E[X_n\mid X_n>0]\longrightarrow\frac m{m-1}.
\end{equation}
Moreover,
\begin{equation}\label{eq:tv}
 \sum_{k\ge1}\left|\Pp(X_n=k\mid X_n>0)-(m-1)m^{-k}\right|
 =O\left(\frac{\log(n+2)}{n+1}\right).
\end{equation}
\end{theorem}

The inclusion of $x=0$ in \eqref{eq:mainprofile} gives
\[
 n^2\rho_n(0)\longrightarrow4.
\]
We explain how the boundary limit $n^2\rho_n(0)\to4$ leads to the survival asymptotic $n^2p_n\to4/(m-1)^2$. By \eqref{eq:notation}--\eqref{eq:rho},
\[
 p_n=\frac{G_n}{m-1}\sum_{j\ge0}m^{-j-1}\rho_n(j).
\]
The estimate on $|\rho_n(j+1)-\rho_n(j)|$ used in Section~\ref{sec:mainproof} shows that replacing $\rho_n(j)$ by $\rho_n(0)$ in this expression introduces an error $o(n^{-2})$. Since $\sum_{j\ge0}m^{-j-1}=1/(m-1)$, we obtain
\[
 p_n=\frac{G_n\rho_n(0)}{(m-1)^2}+o(n^{-2}).
\]
Together with $G_n\to1$ and $n^2\rho_n(0)\to4$, this gives
\[
 n^2p_n\longrightarrow\frac4{(m-1)^2},
\]
which is the second limit in \eqref{eq:main1}.

Under the assumptions of Theorem~\ref{thm:profile}, these conclusions establish Conjectures~1 and~2, the first-moment assertion of Conjecture~3, and Conjecture~4(b) in \cite{CDHLS}. They concern, respectively, the survival probability, the conditional geometric law, the mean, and the generating function at $m$.

The bound \eqref{eq:tv} also gives a rate for convergence of the conditional law. Its left-hand side is twice the total variation distance; this rate does not assert a rate for the convergence of $n^2p_n$. All constants below may depend on $m$ and the fixed initial law, and all limits are taken as $n\to\infty$ unless stated otherwise. The symbol $C$ without a subscript denotes a finite positive constant whose value may change from one occurrence to the next.

\begin{remark}[Sharpness of the moment assumption]
For every $r\in[0,3)$, there exists a nondegenerate critical initial law satisfying $\E[X_0^rm^{X_0}]<\infty$ but $\E[X_0^3m^{X_0}]=\infty$ for which $n^2p_n$ does not converge to $4/(m-1)^2$. Appendix~\ref{sec:sharpness} proves this using a critical law from the family studied in \cite{CDDHLS,CS} and the product estimate of Chen and Shi~\cite[Theorem~1.1]{CS}; see Theorem~\ref{thm:sharpness}. Thus replacing $X_0^3$ by $X_0^r$ for any fixed $0\le r<3$ in the moment assumption does not, in general, guarantee this survival asymptotic.
\end{remark}

\subsection{Proof outline}
We work with the continuous functions $u_N$ defined in \eqref{eq:uN}, which satisfy $u_N(n/N,j/N)=N^2\rho_n(j)$. We first show that every subsequence of $(u_N)_{N\ge1}$ has a further subsequence converging locally uniformly. To bound their values, we combine the generating-function and product estimates in \cite[Theorem~3 and Proposition~1]{CDHLS} with the cubic identity in \cite[eqs.~(18)--(19)]{CDHLS}, and iterate the recurrence relation for $\rho_n$. This gives
\[
 \rho_n(j)\le C\frac{n+1}{(n+j+1)^3},
\]
which controls the size of the rescaled functions $u_N(t,x)$ and their decay as $x\to\infty$.

We next bound the differences of $\rho_n(j)$ as $j$ or $n$ varies. To avoid the factor $\log n$ that would result from adding bounds on $|\rho_n(j+1)-\rho_n(j)|$, we split the convolution sum at a generation chosen according to the increment in $j$. We estimate the earlier and later terms separately. Together with the bounds for differences in $n$, this gives a common modulus of continuity for $u_N$ of order $r(1+\log(1/r))$ on compact rectangles. We can therefore extract locally uniform limits $u$ on $(0,\infty)\times[0,\infty)$, including at $x=0$.

To determine $u$, we need an additional bound on its third moment far from zero. We construct a Markov chain whose law at step $n$ is obtained by multiplying $q_{n,k}$ by the nonnegative polynomial in the cubic identity and normalizing. A bound on its expected positive increments gives
\[
 \int_\delta^\infty x^3u(t,x)\dd x\le C\frac{t^3}{\delta},
 \qquad t,\delta>0.
\]
This estimate controls the remainder in the Laplace transform as $t\downarrow0$ under our third-moment assumption.

To show that all limits coincide, we express the Laplace transform and moment equations in terms of a \emph{string of entrance type}~\cite{Kotani}. The tail estimate determines its spectral measure. Kotani's uniqueness theorem~\cite[Theorem~1]{Kotani} determines the string up to translation. Fixing its right endpoint and using the behavior as $t\downarrow0$ to determine the integration constant gives
\[
 u(t,x)=\frac4{t^2}e^{-2x/t}.
\]
Thus the whole sequence $u_N$ converges locally uniformly to this function. Setting $t=1$ gives
\[
 \sup_{0\le x\le R}\left|n^2\rho_n(\floor{nx})-4e^{-2x}\right|
 \longrightarrow0,
\]
which is \eqref{eq:mainprofile}. The boundary limit and the bounds on $|\rho_n(j+1)-\rho_n(j)|$ give the survival asymptotic, the fixed-$k$ probabilities and the conditional geometric approximation. Summation and the recurrence relation for $G_n$ give the moment and generating-function limits.

For sharpness, Appendix~\ref{sec:sharpness} uses a critical law from the family studied in \cite{CDDHLS,CS}, with $\Pp(X_0=k)\sim c m^{-k}k^{-\alpha}$, $c>0$ and $3<\alpha<4$. The product estimate in \cite[Theorem~1.1]{CS} gives $\prod_{j=0}^{n-1}G_j^{m-1}\asymp n^{\alpha-2}$, whereas criticality and the survival asymptotic $n^2p_n\to4/(m-1)^2$ would imply growth $n^{2+o(1)}$ by Lemma~\ref{lem:productasymptotic}. This contradiction proves Theorem~\ref{thm:sharpness}.

\subsection{Further related work}
A related question is how the free energy and survival probability behave as the initial law approaches criticality. Hu and Shi~\cite{HS} studied free-energy exponents for a generalized Derrida--Retaux model. Chen, Dagard, Derrida, Hu, Lifshits and Shi~\cite{CDDHLS} obtained asymptotics for the free energy in the nearly supercritical case, while Chen, Hu and Shi~\cite{CHSdual} studied the decay of the survival probability and mean in the nearly subcritical case. Chen~\cite{ChenSmooth} also studied infinite differentiability of the free energy.

Some related models can be solved explicitly for particular initial laws. Hu, Mallein and Pain~\cite{HMP} introduced a continuous-time Derrida--Retaux model and studied an exactly solvable family, including its critical scaling. For a model with a geometrically distributed number of children, Li and Zhang~\cite{LZ} obtained asymptotic expansions for geometric-type distributions. Alsmeyer, Hu and Mallein~\cite{AHM} studied the critical curve and free energy for the model on a geometric Galton--Watson tree, including the cases of geometric and exponential initial laws. In our setting, each vertex has exactly $m$ children, and the initial law satisfies the criticality and moment conditions in \eqref{eq:assumptions}.

Scaling limits also connect discrete-time and continuous-time versions of the model. Li and Zhang~\cite{LZscale} proved such limits for generalized Derrida--Retaux models, with the renewal rate varying with the scale parameter. In the limit considered here, the model parameters and initial law remain fixed, and the number of generations tends to infinity.

Derrida, Duquesne and Shi~\cite{DDS} described the continuous critical tree as a growth-fragmentation process. Duquesne and Shi~\cite{DuS} gave a representation in terms of the Brownian continuum random tree.

\subsection{Organization}
Section~\ref{sec:prelim} collects and proves preliminary results on recursion relations. Section~\ref{sec:mainproof} proves Theorem~\ref{thm:profile} using two propositions: one gives bounds on the function $\rho_n(j)$, and one determines every subsequential limit of $\rho_{n}(j)$ after scaling. The two propositions are proved in Sections~\ref{sec:pointwise}--\ref{sec:rigidity}. Appendix~\ref{sec:sharpness} proves Theorem~\ref{thm:sharpness} for the sharpness of the moment condition \eqref{eq:assumptions}.

\subsection{AI-assisted proof development and verification status}
\label{subsec:ai-declaration}

Generative AI was used to produce the mathematical proof draft presented
in this manuscript. OpenAI ChatGPT and Codex also assisted with manuscript
preparation and organization. The authors have manually checked the
mathematical arguments in the current draft and take responsibility for
its content.

A Lean~4 formalization of an earlier manuscript snapshot is available in
the accompanying \href{https://github.com/iyalice/DR_Yaglom_lean/tree/2830c58dee0265b8cc71b0dc16acc944cda316a2/lean}{Lean repository}.
The repository records a successful Lean~4.19.0 build and axiom audit on
12 September 2026. The formalized results include the profile theorem
and its asymptotic consequences, as stated in Theorem~\ref{thm:profile}.
The development accepts four explicitly declared external inputs: two
estimates from~\cite{CDHLS}, the inverse uniqueness input from~\cite{Kotani},
and the stable-product estimate from~\cite{CS}. These inputs are not
re-proved in the Lean project; their signatures and uses are documented
there.

The formalization is tied to the frozen manuscript supplied with that
repository. In particular, its sharpness theorem concerns the failure of
$n(G_n-1)\to2/(m-1)$, whereas Theorem~\ref{thm:sharpness} in the present
revision concerns the failure of $n^2p_n\to4/(m-1)^2$. The latter statement
is not covered by the existing formalized sharpness theorem. A complete
correspondence check against the present revision remains to be carried
out; formal verification of the entire current manuscript is therefore
not claimed.

\section{Preliminaries}\label{sec:prelim}
Throughout Sections~\ref{sec:prelim}--\ref{sec:rigidity}, we assume the hypotheses of Theorem~\ref{thm:profile}. We collect the tilted recursion and moment estimates used in the proof of Theorem~\ref{thm:profile}.

\subsection{Exact tilted recursion and coarse bounds}
Let $\Gamma_n(s)=\E[s^{X_n}]$, and recall that $G_n=\Gamma_n(m)$ and $\varepsilon_n=G_n-1$. For the tilted law in \eqref{eq:notation}, put
\[
 H_n(z)=\sum_{k\ge0}q_{n,k}z^k,\quad x_n=q_{n,0},\quad
 \theta_n=1-x_n,\quad D_n=1+(m-1)x_n^m.
\]

\begin{proposition}\label{prop:pgf}
For every $n\ge0$, $\E[X_n^3m^{X_n}]<\infty$, and, for $0<s\le m$,
\begin{equation}\label{eq:pgf}
 \Gamma_{n+1}(s)=\frac{\Gamma_n(s)^m-\Gamma_n(0)^m}{s}+\Gamma_n(0)^m.
\end{equation}
Criticality is preserved:
\[
 m(m-1)\Gamma'_n(m)=\Gamma_n(m).
\]
Equivalently, the tilted law satisfies
\begin{equation}\label{eq:qmean}
 \sum_{k\ge0}q_{n,k}=1,\qquad \sum_{k\ge0}kq_{n,k}=\frac1{m-1}.
\end{equation}
Consequently, the generating function $H_n$ of $q_n$ and the values $G_n=\Gamma_n(m)$ satisfy, for $0<z\le1$,
\begin{equation}\label{eq:tilt}
 H_{n+1}(z)=\frac{H_n(z)^m+(mz-1)x_n^m}{zD_n},\qquad
 G_{n+1}=\frac{G_n^mD_n}{m}.
\end{equation}
\end{proposition}
\begin{proof}
Finiteness of $\E[X_n^3m^{X_n}]$, the recursion \eqref{eq:pgf}, and preservation of criticality follow from \cite[Section~2.1, eqs.~(7)--(13)]{CDHLS}. Dividing the criticality identity by $G_n$ gives the mean in \eqref{eq:qmean}; the normalization follows from the definition of $q_n$.

To obtain \eqref{eq:tilt}, use $\Gamma_n(0)=G_nx_n$ in \eqref{eq:pgf} at $s=m$, which gives
\[
 G_{n+1}=\frac{G_n^m}{m}\bigl(1+(m-1)x_n^m\bigr)=\frac{G_n^mD_n}{m}.
\]
Substituting $s=mz$ and dividing by $G_{n+1}$ then yields the recursion for $H_n(z)=\Gamma_n(mz)/G_n$.
\end{proof}

\begin{proposition}\label{prop:coarse}
There is a constant $C<\infty$, depending on $m$ and the initial law, such that, for all $n\ge0$,
\begin{equation}\label{eq:coarse}
 0\le\varepsilon_n\le\frac C{n+1},\qquad
 \prod_{j=0}^{n-1}G_j^{m-1}\le C(n+1)^2.
\end{equation}
Moreover,
\begin{equation}\label{eq:theta}
 p_n\le\frac{\varepsilon_n}{m-1},\qquad
 \theta_n=\frac{\varepsilon_n+p_n}{G_n}=O((n+1)^{-1}).
\end{equation}
In particular, $x_n\to1$ and $\inf_{n\ge0}x_n>0$.
\end{proposition}
\begin{proof}
By \cite[Theorem~3 and Proposition~1]{CDHLS}, for $n\ge1$,
\[
 G_n-1\le c_2\frac{\E[X_0^3m^{X_0}]^{1/2}}{\Pp(X_0=0)^{1/2}}\frac1n,
 \qquad
 \prod_{j=0}^{n}G_j^{m-1}\le\frac{c_3}{\Pp(X_0=0)}n^2,
\]
where $c_2,c_3>0$ depend only on $m$. Since criticality and the exclusion of the binary fixed point imply $\Pp(X_0=0)>0$, these estimates yield \eqref{eq:coarse}.

Since $m^k-1\ge m-1$ for $k\ge1$,
\[
 \varepsilon_n=\sum_{k\ge1}(m^k-1)\Pp(X_n=k)
 \ge(m-1)\sum_{k\ge1}\Pp(X_n=k)=(m-1)p_n.
\]
This proves the first inequality in \eqref{eq:theta}. Also,
\[
 \theta_n=\sum_{k\ge1}q_{n,k}
 =\frac{G_n-\Pp(X_n=0)}{G_n}
 =\frac{\varepsilon_n+p_n}{G_n}.
\]
Since $G_n\ge1$, the preceding inequality and \eqref{eq:coarse} give
$\theta_n\le m\varepsilon_n/(m-1)=O((n+1)^{-1})$, hence $x_n\to1$.
Finally, \eqref{eq:model} implies
$\Pp(X_{n+1}=0)\ge\Pp(X_n=0)^m>0$ by induction, so $x_n>0$ for every $n$. Together with $x_n\to1$, this yields $\inf_{n\ge0}x_n>0$.
\end{proof}

\subsection{Moment bounds}
The next identity is the normalized form of the cubic generating-function identity in \cite[eqs.~(18)--(19)]{CDHLS}. Set $\chi_m=1-1/(m-1)$ and
\begin{equation}\label{eq:cubic}
 h_m(k)=k(k+1)\left(k-\frac1{m-1}\right)
 =(k)_3+(3+\chi_m)(k)_2+2\chi_m k,
\end{equation}
where $(k)_r=k(k-1)\cdots(k-r+1)$.

\begin{lemma}\label{lem:cubic}
Let $K$ be a random variable with law $q_n$, and let $J_n=\E h_m(K) = \sum^{\infty}_{k = 0} h_m(k) q_{n, k}$. Then $J_n>0$ and
\begin{equation}\label{eq:J}
 D_nJ_{n+1}=mJ_n,\qquad
 J_n=J_0\frac{G_0}{G_n}\prod_{j=0}^{n-1}G_j^{m-1}\le C(n+1)^2.
\end{equation}
Consequently,
\begin{equation}\label{eq:lowmoments}
 \E K^3\le C(n+1)^2,\qquad
 \E K^2\le C(n+1).
\end{equation}
\end{lemma}
\begin{proof}
We write $\Xi_n$ for the quantity $\mathcal D_n(m)$ defined in \cite[eq.~(18)]{CDHLS}; explicitly,
\[
 \Xi_n:=m(m-1)\Gamma_n'''(m)
 +(4m-5)\Gamma_n''(m)
 +\frac{2(m-2)}{m^2(m-1)}G_n.
\]
Using \eqref{eq:cubic}, $\E[(K)_r]=m^r\Gamma_n^{(r)}(m)/G_n$, and $\E K=1/(m-1)$, we obtain
\[
\begin{aligned}
 J_n
 &=\frac{m^3}{G_n}\Gamma_n'''(m)
   +\frac{(4m-5)m^2}{(m-1)G_n}\Gamma_n''(m)
   +\frac{2(m-2)}{(m-1)^2}\\
 &=\frac{m^2}{(m-1)G_n}
   \left[
   m(m-1)\Gamma_n'''(m)
   +(4m-5)\Gamma_n''(m)
   +\frac{2(m-2)}{m^2(m-1)}G_n
   \right]\\
 &=\frac{m^2}{(m-1)G_n}\Xi_n.
\end{aligned}
\]
The identity $\Xi_{n+1}=\Xi_nG_n^{m-1}$ from \cite[eq.~(19)]{CDHLS} therefore gives
\[
 J_{n+1}=\frac{G_n^m}{G_{n+1}}J_n=\frac{m}{D_n}J_n.
\]
Iteration yields the product formula in \eqref{eq:J}, and \eqref{eq:coarse} gives its upper bound. For $m\ge3$, $J_0>0$ follows from $h_m(k)>0$ for $k\ge1$ and \eqref{eq:qmean}. In the binary case, criticality and nondegeneracy force positive mass above one, so $J_0>0$ as well. The recursion then gives $J_n>0$ for every $n$.

For $m\ge3$, $h_m(k)\ge k^3$ for integers $k\ge1$. For $m=2$, $k^3\le(4/3)h_2(k)$ when $k\ge2$, and the contribution at one is bounded. Thus the cubic bound in \eqref{eq:lowmoments} follows from \eqref{eq:coarse}. Cauchy--Schwarz gives $\E K^2\le(\E K\,\E K^3)^{1/2}$, proving the second bound.
\end{proof}

\subsection{Tilted recursion and auxiliary estimates}\label{subsec:transport}
In this subsection, we first express the tilted recursion in terms of probability laws, then identify
its linear coefficient and derive the estimates used below.

\begin{lemma}[Distributional recursion for the tilted laws]\label{lem:tiltedlaw}
For each $n\ge0$, let $K_{n,1},\ldots,K_{n,m}$
be independent with law $q_n$, and put $S_n=\sum_{i=1}^mK_{n,i}$. Then
\[
 q_{n+1}
 =\frac1{D_n}\operatorname{Law}\bigl((S_n-1)_+\bigr)
  +\frac{(m-1)x_n^m}{D_n}\delta_0,
\]
where $\delta_0$ is the point mass at zero.
\end{lemma}
\begin{proof}
For inputs summing to $s$, the product of their exponential weights is $m^s$,
whereas the output $(s-1)_+$ has weight $m^{(s-1)_+}$. The ratio of the latter
to the former is $1/m$ for $s\ge1$ and $1$ for $s=0$. After removing the
common factor $1/m$, all inputs retain weight one except the all-zero input,
which has weight $m$. This adds mass $(m-1)x_n^m$ at zero to the law of
$(S_n-1)_+$. Dividing by the total mass $D_n=1+(m-1)x_n^m$ gives the result.
\end{proof}

By Lemma~\ref{lem:tiltedlaw}, $q_{n+1,k}=\Pp(S_n=k+1)/D_n$ for $k\ge1$.
Splitting this probability according to the number of positive inputs gives
\[
 q_{n+1,k}
 =\frac{mx_n^{m-1}}{D_n}q_{n,k+1}
 +\frac1{D_n}\sum_{r=2}^m\binom mr x_n^{m-r}
 \sum_{\substack{k_1,\ldots,k_r\ge1\\k_1+\cdots+k_r=k+1}}
 \prod_{i=1}^r q_{n,k_i}.
\]
The first term is linear in the positive masses and corresponds to exactly one positive input; the remaining terms involve at least two.
We denote the linear coefficient by $c_n$ and the product of these coefficients by $C_n$:
\begin{equation}\label{eq:c}
 c_n=\frac{mx_n^{m-1}}{D_n},\qquad C_0=1,\qquad C_n=\prod_{j=0}^{n-1}c_j.
\end{equation}
\begin{lemma}\label{lem:transportcoefficients}
The coefficients defined in \eqref{eq:c} satisfy
\begin{equation}\label{eq:C}
 0<c_n\le1,\qquad 1-c_n=O((n+1)^{-2}),\qquad
 C_n\downarrow C_\infty>0.
\end{equation}
\end{lemma}
\begin{proof}
Proposition~\ref{prop:coarse} gives $x_n>0$ and
$\theta_n=1-x_n=O((n+1)^{-1})$. By \eqref{eq:c}, $c_n>0$ and
\[
 1-c_n=\frac{1+(m-1)x_n^m-mx_n^{m-1}}{D_n}.
\]
The numerator is nonnegative on $[0,1]$ and vanishes to second order at $1$.
Since $D_n\ge1$, it follows that
\[
 0\le1-c_n\le C\theta_n^2=O((n+1)^{-2}).
\]
Thus $0<c_n\le1$. Finally, $\sum_{n\ge0}(1-c_n)<\infty$ ensures that
the product of the positive factors $c_n$ has a positive limit,
so $C_n\downarrow C_\infty>0$.
\end{proof}

For the mass at zero, Lemma~\ref{lem:tiltedlaw} gives
\[
 x_{n+1}=\frac{mx_n^m+mx_n^{m-1}q_{n,1}}{D_n}
 =c_n(x_n+q_{n,1}).
\]
We will use this recursion for $x_n=q_{n,0}$ to estimate the sum in the next lemma, which will be needed in the proof of Proposition~\ref{prop:far}.
\begin{lemma}\label{lem:atomtail}
The sequence $(q_{n,1})_{n\ge0}$ is summable, and its tails satisfy
\begin{equation}\label{eq:atomtail}
 T_n:=\sum_{s\ge n}q_{s,1}\le\frac C{n+1}.
\end{equation}
\end{lemma}
\begin{proof}
Rearranging the recursion for $x_n$ gives
\begin{equation}\label{eq:atomidentity}
 q_{n,1}=\theta_n-\theta_{n+1}+(c_n^{-1}-1)x_{n+1}.
\end{equation}
By Lemma~\ref{lem:transportcoefficients}, $\inf_s c_s>0$ and
$1-c_s=O((s+1)^{-2})$, so
$0\le(c_s^{-1}-1)x_{s+1}\le C(s+1)^{-2}$.
Summing \eqref{eq:atomidentity} over $n\le s\le r$ therefore gives
\[
 \sum_{s=n}^r q_{s,1}
 =\theta_n-\theta_{r+1}+\sum_{s=n}^r(c_s^{-1}-1)x_{s+1}
 \le\theta_n+C\sum_{s=n}^r(s+1)^{-2}.
\]
Letting $r\to\infty$ and using $\theta_n=O((n+1)^{-1})$
from Proposition~\ref{prop:coarse}, we obtain
\[
 T_n\le\theta_n+C\sum_{s=n}^{\infty}(s+1)^{-2}
 \le\frac C{n+1},
\]
as claimed.
\end{proof}

We conclude this section by collecting the moment estimates for $\rho_n$ in a weighted norm.
Put $L_n=n+1$. For a sequence $f$ on $\Zp$ and $L\ge1$, define
\[
 \norm f_{1,3,L}=\sum_{j\ge0}(1+j/L)^3|f(j)|.
\]
Recall that $\rho_n(j)=(m-1)q_{n,j+1}$ for $j\in\Zp$.
\begin{lemma}\label{lem:weightedmass}
There exists a constant $C>0$, depending only on $m$ and the law of $X_0$, such that for every $n\ge0$,
\begin{equation}\label{eq:weightedmass}
 \norm{\rho_n}_{1,3,L_n}\le\frac C{L_n},\qquad
 \sum_{j\ge0}(j+1)\rho_n(j)=1.
\end{equation}
\end{lemma}
\begin{proof}
By \eqref{eq:qmean}, \eqref{eq:theta}, and \eqref{eq:lowmoments},
\[
 \sum_{j\ge0}\rho_n(j)=(m-1)\theta_n\le\frac C{L_n},\qquad
 \sum_{j\ge0}j^r\rho_n(j)\le CL_n^{r-1},\quad r=1,2,3.
\]
Expanding $(1+j/L_n)^3$ gives the first bound. The identity follows from
\[
 \sum_{j\ge0}(j+1)\rho_n(j)
 =(m-1)\sum_{k\ge1}kq_{n,k}=1.\qedhere
\]
\end{proof}

\section{Proof of Theorem~\ref{thm:profile}}\label{sec:mainproof}
In this section, we prove Theorem~\ref{thm:profile}.
We first state the two propositions used in the proof. The first proposition bounds $\rho_n$ and its differences in space and time; these bounds will allow us to extract locally uniformly convergent subsequences of the rescaled functions $u_N$ defined below. The second proposition determines all such limits of $u_N$. The proofs of these propositions are given in Sections~\ref{sec:pointwise}--\ref{sec:rigidity} and use the assumptions of Theorem~\ref{thm:profile} and the estimates in Section~\ref{sec:prelim}.

For a sequence $f$ on $\Zp$ and $L\ge1$, define
\[
 \norm{f}_{\infty,3,L}
 :=\sup_{j\ge0}(1+j/L)^3|f(j)|.
\]
Define the operators $S$ and $\Delta$ by
\[
 (Sf)(j)=f(j+1),\qquad
 (\Delta f)(j)=f(j)-f(j+1),\qquad j\ge0.
\]
For each integer $h\ge 1$, $S^h$ denotes the $h$-fold
composition of $S$; thus
$(S^h f)(j)=f(j+h)$.

\begin{proposition}\label{prop:smoothing}
Under the assumptions of Theorem~\ref{thm:profile}, there exists a constant $C>0$, depending only on $m$ and the law of $X_0$, such that the following bounds hold for every integer $n\ge0$:
\begin{align}
 \norm{\rho_n}_{\infty,3,n+1}&\le C(n+1)^{-2},\label{eq:sup}\\
 \norm{\Delta\rho_n}_{\infty,3,n+1}&\le C\frac{\log(n+2)}{(n+1)^3}.\label{eq:gradient}
\end{align}
For every integer $1\le h\le n+1$,
\begin{equation}\label{eq:spmod}
 \norm{\rho_n-S^h\rho_n}_{\infty,3,n+1}
 \le \frac{Ch}{(n+1)^3}\left(1+\log\frac{n+1}{h}\right).
\end{equation}
The differences in time satisfy the same bound in the ordinary supremum norm:
\begin{equation}\label{eq:timemod}
 \norm{\rho_{n+h}-\rho_n}_{\infty}
 \le \frac{Ch}{(n+1)^3}\left(1+\log\frac{n+1}{h}\right).
\end{equation}
\end{proposition}

To state the second proposition, we first define the rescaled function $u_N: (0,\infty)\times
[0,\infty) \to \mathbb{R}$. For each integer $N\ge1$, set
\begin{equation}\label{eq:rescaled}
 u_N(n/N,j/N) := N^2\rho_n(j),\qquad n,j\ge0.
\end{equation}
On each small rectangle $[n/N,(n+1)/N]\times[j/N,(j+1)/N]$, define $u_N$ by bilinear interpolation: for $\alpha,\beta\in[0,1]$, set
\begin{equation}\label{eq:uN}
 \begin{split}
 u_N\left(\frac{n+\alpha}{N},\frac{j+\beta}{N}\right)
 &:=N^2(1-\alpha)\bigl[(1-\beta)\rho_n(j)+\beta\rho_n(j+1)\bigr]\\
 &\quad+N^2\alpha\bigl[(1-\beta)\rho_{n+1}(j)+\beta\rho_{n+1}(j+1)\bigr].
 \end{split}
\end{equation}
Then $u_N$ is continuous on $(0,\infty)\times
[0,\infty)$. 
The next proposition determines every locally uniform subsequential limit of $u_N$ as $N\to\infty$. Once we have proved compactness for $(u_N)_{N\ge 1}$, this will give convergence of the whole sequence. We prove it in Section~\ref{sec:rigidity}, using the moment bounds and limit equations from Section~\ref{sec:limits}.

\begin{proposition}\label{prop:identification}
Under the assumptions of Theorem~\ref{thm:profile}, every locally uniform subsequential limit of $u_N$ on $(0,\infty)\times[0,\infty)$ is
\begin{equation}\label{eq:identifiedprofile}
 u(t,x)=\frac4{t^2}e^{-2x/t}.
\end{equation}
\end{proposition}

We are now ready to prove Theorem~\ref{thm:profile}.
\begin{proof}[Proof of Theorem~\ref{thm:profile}]
We first prove the convergence stated in \eqref{eq:mainprofile}.

We begin by proving that $(u_N)_{N\ge1}$ is uniformly bounded and equicontinuous on each compact rectangle in $(0,\infty)\times[0,\infty)$. Arzel\`a--Ascoli will then give subsequences that converge uniformly on each such rectangle. Fix a compact rectangle $K=[\eta,T]\times[0,R]$, where $0<\eta<T<\infty$ and $R<\infty$. For all sufficiently large $N$, every vertex $(n/N,j/N)$ of a small rectangle intersecting $K$ satisfies $n/N\in[\eta/2,T+1]$. At these points, \eqref{eq:sup} gives
\[
 |u_N(n/N,j/N)|=N^2\rho_n(j)
 \le C\frac{N^2}{(n+1)^2}\le \frac{4C}{\eta^2}=: C_\eta.
\]
By \eqref{eq:uN}, each value of $u_N$ in a small rectangle is a convex combination of the four vertex values. Thus $|u_N(t,x)|\le C_\eta$ for every $(t,x)\in K$ and all sufficiently large $N$.

To prove equicontinuity of $(u_N)_{N\ge 1}$ on $K$, we first bound the differences of $u_N(t,x)$ between grid points that differ only in $x$ or only in $t$. We will then use \eqref{eq:uN} to extend these bounds to all points of $K$. Put $\omega(0):=0$ and $\omega(r):=r(1+\log(1/r))$ for $0<r\le1$. For grid times $n/N \in [\eta/2,T+1]$ and $1\le h\le\min\{n+1,N\}$, equations \eqref{eq:spmod} and \eqref{eq:timemod} imply
\begin{equation}\label{eqn:diff grid points}
  \begin{split}
 N^2\bigl(\norm{\rho_n-S^h\rho_n}_\infty
       +\norm{\rho_{n+h}-\rho_n}_\infty\bigr)
  &\le \frac{2CN^2h}{(n+1)^3}
       \left(1+\log\frac{n+1}{h}\right)\\
 &\le C_{\eta,T}\,\omega(h/N).
 \end{split}
\end{equation}
The last bound uses $\eta/2\le(n+1)/N\le T+2$. Here $C_{\eta,T}$ is independent of $N,n,h$.
For grid points $z=(t,x)$ and $z'=(t',x')$ with $t,t'\in[\eta/2,T+1]$, let $r:=\max\{|t-t'|,|x-x'|\}\le\min\{1,\eta/2\}$. Applying \eqref{eqn:diff grid points} in each coordinate direction and using the monotonicity of $\omega$ gives
\[
 \begin{split}
 |u_N(z)-u_N(z')|
 &\le |u_N(t,x)-u_N(t,x')|
       +|u_N(t,x')-u_N(t',x')|\\
 &\le 2C_{\eta,T}\omega(r).
 \end{split}
\]

To extend the estimate to arbitrary points, for $z=(t,x)\in K$, set $z_N:=(\floor{Nt}/N,\floor{Nx}/N)$. By \eqref{eq:uN}, $u_N(z)$ is a convex combination of the values at the four vertices of the small rectangle containing $z$. Each vertex can be reached from $z_N$ along at most two sides of this rectangle, so \eqref{eqn:diff grid points} with $h=1$ gives
\begin{equation}\label{eqn:diff_grid_vs_nongird}
   |u_N(z)-u_N(z_N)|
 \le 2C_{\eta,T}\omega(1/N)
 =2C_{\eta,T}\frac{1+\log N}{N}.
\end{equation}
Now let $z=(t,x),z'=(t',x')\in K$, with corresponding grid points $z_N,z_N'$, and let $r := \max\{|t-t'|,|x-x'|\}$. If $N^{-1}\le r\le\min\{1/3,\eta/6\}$, then $z_N$ and $z_N'$ differ by at most $r+2/N\le3r$ in each coordinate. Combining these estimates with \eqref{eqn:diff_grid_vs_nongird} gives
\[
 \begin{split}
 |u_N(z)-u_N(z')|
 &\le |u_N(z)-u_N(z_N)|+|u_N(z_N)-u_N(z_N')|\\
 &\quad+|u_N(z_N')-u_N(z')|\\
 &\le 2C_{\eta,T}\omega(3r)+4C_{\eta,T}\omega(1/N)\\
 &\le 10C_{\eta,T}\omega(r).
 \end{split}
\]
The last step uses $\omega(3r)\le3\omega(r)$ and
$\omega(1/N)\le\omega(r)$.

For $0<r<1/N$, first compare points in the same small rectangle that differ in one coordinate. Write $t=(n+\alpha)/N$, with $0\le\alpha\le1$, and take $x,x'\in[j/N,(j+1)/N]$. Subtracting the two expressions in \eqref{eq:uN} gives
\[
 \begin{split}
 u_N(t,x)-u_N(t,x')
 ={}&N^3(x-x')\bigl[(1-\alpha)(\rho_n(j+1)-\rho_n(j))\\
 &\qquad\qquad+\alpha(\rho_{n+1}(j+1)-\rho_{n+1}(j))\bigr].
 \end{split}
\]
Applying \eqref{eqn:diff grid points} with $h=1$, we have 
\begin{equation}\label{eq:space diff}
  |u_N(t,x)-u_N(t,x')|\le C_{\eta,T}(1+\log N)|x-x'|.
\end{equation}
The same calculation in the time coordinate gives, for $t,t'\in[n/N,(n+1)/N]$ in the same small rectangle,
\begin{equation}\label{eq:time diff}
 |u_N(t,x)-u_N(t',x)|\le C_{\eta,T}(1+\log N)|t-t'|.
\end{equation}
For arbitrary $z,z'\in K$ with $0<r<1/N$, compare values along $(t,x)\to(t,x')\to(t',x')$. Since each segment has length less than $1/N$, it can be divided into at most two parts, each contained in one small rectangle. Apply \eqref{eq:space diff} to each horizontal part and \eqref{eq:time diff} to each vertical part. The lengths add to $|x-x'|$ and $|t-t'|$, respectively, so the triangle inequality gives \eqref{eq:space diff} and \eqref{eq:time diff} for $x, x'$ and $t,t'$ not in the same small rectangle. Hence,
\[
 \begin{split}
 |u_N(z)-u_N(z')|
 &\le |u_N(t,x)-u_N(t,x')|+|u_N(t,x')-u_N(t',x')|\\
 &\le C_{\eta,T}(1+\log N)\bigl(|x-x'|+|t-t'|\bigr)\\
 &\le 2C_{\eta,T}r(1+\log N)\\
 &\le 2C_{\eta,T}r\left(1+\log\frac1r\right)
 =2C_{\eta,T}\omega(r).
 \end{split}
\]

Thus $(u_N)_{N\ge1}$ is equicontinuous on $K$.

Let $(u_{N_j})_{j\ge1}$ be any subsequence of $(u_N)_{N\ge1}$. We will extract a further subsequence that converges locally uniformly on $(0,\infty)\times[0,\infty)$. Apply Arzel\`a--Ascoli to $(u_{N_j})_{j\ge1}$ on the nested compact sets
\[
 K_\ell=[1/\ell,\ell]\times[0,\ell],\qquad \ell\ge2.
\]
At each step choose a subsequence of the previous one. More precisely, let
$I_\ell$ be the infinite set of indices retained at step $\ell$, so that
$I_{\ell+1}\subset I_\ell$ and $(u_{N_j})_{j\in I_\ell}$ converges uniformly
on $K_\ell$. Choose $j_2\in I_2$, and then choose $j_k\in I_k$ recursively
with $j_k>j_{k-1}$ for $k\ge3$. For every fixed $\ell$, we have $j_k\in I_\ell$ whenever
$k\ge\ell$. Thus the single diagonal subsequence $(u_{N_{j_k}})_{k\ge2}$
converges uniformly on every $K_\ell$. Its limits agree on overlaps and
define a function $v$ on $(0,\infty)\times[0,\infty)$. Every compact subset
of this domain lies in some $K_\ell$, so the convergence is locally uniform,
including at $x=0$.

Proposition~\ref{prop:identification} shows that $v$ equals the function $u$ in \eqref{eq:identifiedprofile}. Since the original subsequence $(u_{N_j})_{j\ge1}$ was arbitrary, every subsequence has a further subsequence converging locally uniformly to $u$. If the whole sequence did not converge locally uniformly to $u$, there would be a compact set $K$, a number $\epsilon>0$, and a subsequence such that
\[
 \sup_{z\in K}|u_{N_j}(z)-u(z)|\ge\epsilon\qquad\text{for every }j.
\]
A further subsequence converging uniformly to $u$ on $K$ contradicts this inequality. Hence the whole sequence $(u_N)_{N\ge1}$ converges locally uniformly to
\[ u(t,x)=\frac4{t^2}e^{-2x/t}.\]

We now show that $n^2\rho_n(\floor{nx})$ converges to $4e^{-2x}$ uniformly on $[0,R]$. Set $N=n$ and $t=1$. Equations \eqref{eq:uN} and \eqref{eq:gradient} give
\[
 \sup_{0\le x\le R}
 \left|u_n(1,x)-n^2\rho_n(\floor{nx})\right|
 \le n^2\norm{\Delta\rho_n}_\infty
 =O\left(\frac{\log(n+2)}{n+1}\right)\longrightarrow0.
\]
Combining this estimate with the uniform convergence of $u_n(1,x)$ on $[0,R]$, we obtain
\[
 \begin{split}
 &\sup_{0\le x\le R}\left|n^2\rho_n(\floor{nx})-4e^{-2x}\right|\\
 &\quad\le \sup_{0\le x\le R}\left|n^2\rho_n(\floor{nx})-u_n(1,x)\right|
       +\sup_{0\le x\le R}\left|u_n(1,x)-4e^{-2x}\right|
 \longrightarrow0.
 \end{split}
\]
This is the convergence stated in \eqref{eq:mainprofile}. In particular, taking $x=0$ gives
\begin{equation}\label{eq:boundary4}
 n^2\rho_n(0)\longrightarrow4.
\end{equation}

We next derive $n^2p_n\to4/(m-1)^2$ from \eqref{eq:boundary4}. The definitions of $q_n$ and $\rho_n$ give
\[
 \Pp(X_n=k)=\frac{G_n}{m-1}m^{-k}\rho_n(k-1),
 \qquad k\ge1.
\]
Summing over $k$ gives the survival probability
\begin{equation*}
 p_n=\frac{G_n}{m-1}\sum_{j\ge0}m^{-j-1}\rho_n(j).
\end{equation*}
To replace $\rho_n(j)$ by $\rho_n(0)$ in this sum, use
\[
 |\rho_n(j)-\rho_n(0)|
 \le\sum_{\ell=0}^{j-1}|\rho_n(\ell+1)-\rho_n(\ell)|
 \le j\norm{\Delta\rho_n}_\infty,
\]
where the sum is empty when $j=0$. Since
\[
 \sum_{j\ge0}m^{-j-1}=\frac1{m-1},
 \qquad
 \sum_{j\ge0}jm^{-j-1}=\frac1{(m-1)^2},
\]
we obtain
\[
 \begin{split}
 \left|p_n-\frac{G_n\rho_n(0)}{(m-1)^2}\right|
 &\le\frac{G_n}{m-1}
       \sum_{j\ge0}m^{-j-1}|\rho_n(j)-\rho_n(0)|\\
 &\le\frac{G_n\norm{\Delta\rho_n}_\infty}{(m-1)^3}
 =O\left(\frac{\log(n+2)}{(n+1)^3}\right).
 \end{split}
\]
Multiplying by $n^2$ and using $G_n\to1$ and $n^2\rho_n(0)\to4$, we obtain
\[
 n^2p_n=\frac{G_n}{(m-1)^2}\,n^2\rho_n(0)+o(1)
 \longrightarrow\frac4{(m-1)^2}.
\]
This is the second limit in \eqref{eq:main1}.

To compute the limit of $n(G_n-1)$, we first sum $\rho_n$ over $j$. For fixed $A\ge1$, \eqref{eq:mainprofile} and Riemann sums give
\[
 n\sum_{0\le j\le An}\rho_n(j)
 =\frac1n\sum_{0\le j\le An}n^2\rho_n(j)
 \longrightarrow\int_0^A4e^{-2x}\dd x.
\]
To control the tail, use $\rho_n(j)\le C(n+1)/(n+j+1)^3$ from \eqref{eq:sup} and compare the sum with an integral:
\[
 n\sum_{j>An}\rho_n(j)
 \le Cn(n+1)\sum_{j>An}\frac1{(n+j+1)^3}
 \le\frac{C}{(1+A)^2}.
\]
Letting first $n\to\infty$ and then $A\to\infty$ gives
\[
 n\sum_{j\ge0}\rho_n(j)
 \longrightarrow\int_0^\infty4e^{-2x}\dd x=2.
\]
Since $\sum_j\rho_n(j)=(m-1)\theta_n$, we have $n\theta_n\to2/(m-1)$. Using $\varepsilon_n=G_n\theta_n-p_n$ from \eqref{eq:theta}, together with $G_n\to1$ and $p_n=O(n^{-2})$, we obtain
\[
 n(G_n-1)=n\varepsilon_n=G_n(n\theta_n)-np_n
 \longrightarrow\frac2{m-1}.
\]
This is the first limit in \eqref{eq:main1}.

For the limit of $n^2\Pp(X_n=k)$ with fixed $k\ge1$, compare $\rho_n(k-1)$ with $\rho_n(0)$:
\[
 n^2|\rho_n(k-1)-\rho_n(0)|
 \le(k-1)n^2\norm{\Delta\rho_n}_\infty\longrightarrow0.
\]
Equation~\eqref{eq:boundary4} therefore gives
\[
 n^2\Pp(X_n=k)
 =\frac{G_n}{m-1}m^{-k}\,n^2\rho_n(k-1)
 \longrightarrow\frac4{(m-1)m^k}.
\]
This is the first limit in \eqref{eq:atoms}.
For the mean, \eqref{eq:sup} and boundedness of $G_n$ give $n^2k\Pp(X_n=k)\le Ckm^{-k}$, a summable bound independent of $n$. Dominated convergence gives
\[
 n^2\E[X_n]\longrightarrow\frac4{m-1}\sum_{k\ge1}km^{-k}
 =\frac{4m}{(m-1)^3}.
\]
This is the second limit in \eqref{eq:main2}.
Dividing by the survival probability and using $n^2p_n\to4/(m-1)^2$, we obtain
\[
 \E[X_n\mid X_n>0]
 =\frac{n^2\E[X_n]}{n^2p_n}
 \longrightarrow
 \frac{4m/(m-1)^3}{4/(m-1)^2}
 =\frac m{m-1}.
\]
This is the second limit in \eqref{eq:atoms}.

We next bound the sum of the absolute differences between the conditional probabilities $\Pp(X_n=k\mid X_n>0)$ and $(m-1)m^{-k}$. Put
\[
 a_n=\sum_{k\ge1}m^{-k}\rho_n(k-1)
     =\frac{m-1}{G_n}p_n
     \sim\frac4{(m-1)n^2}.
\]
Then
\[
 \Pp(X_n=k\mid X_n>0)
 =\frac{m^{-k}\rho_n(k-1)}{a_n}.
\]
To compare these probabilities with $(m-1)m^{-k}$, use $|\rho_n(j)-\rho_n(0)|\le j\norm{\Delta\rho_n}_\infty$ to obtain
\[
 \left|a_n-\frac{\rho_n(0)}{m-1}\right|
 \le \norm{\Delta\rho_n}_\infty\sum_{k\ge1}(k-1)m^{-k}
 =\frac{\norm{\Delta\rho_n}_\infty}{(m-1)^2}.
\]
The triangle inequality then gives
\[
 \begin{split}
 |\rho_n(k-1)-(m-1)a_n|
 &\le|\rho_n(k-1)-\rho_n(0)|
       +|\rho_n(0)-(m-1)a_n|\\
 &\le \norm{\Delta\rho_n}_\infty\left(k-1+\frac1{m-1}\right).
 \end{split}
\]
Multiplying by $m^{-k}/a_n$, summing over $k\ge1$, and using \eqref{eq:gradient} gives
\[
 \begin{split}
 &\sum_{k\ge1}
   \left|\Pp(X_n=k\mid X_n>0)-(m-1)m^{-k}\right|\\
 &\qquad\le\frac{\norm{\Delta\rho_n}_\infty}{a_n}
       \sum_{k\ge1}m^{-k}\left(k-1+\frac1{m-1}\right)
  =\frac{2\norm{\Delta\rho_n}_\infty}{(m-1)^2a_n}
  =O\left(\frac{\log(n+2)}{n+1}\right).
  \end{split}
\]
This is the bound stated in \eqref{eq:tv}.

Finally, we compute the limit of $n^2(\varepsilon_n-\varepsilon_{n+1})$ from the generating function recursion. Evaluating \eqref{eq:pgf} at $m$ gives
\[
 G_{n+1}=\frac{G_n^m}{m}+\frac{m-1}{m}(1-p_n)^m.
\]
Writing $G_n=1+\varepsilon_n$ gives
\[
 m(\varepsilon_n-\varepsilon_{n+1})
 =(m-1)\{1-(1-p_n)^m\}-\{(1+\varepsilon_n)^m-1-m\varepsilon_n\}.
\]
The binomial formula gives
\[
 \begin{split}
 1-(1-p_n)^m&=mp_n+O(p_n^2),\\
 (1+\varepsilon_n)^m-1-m\varepsilon_n
 &=\frac{m(m-1)}2\varepsilon_n^2+O(\varepsilon_n^3).
 \end{split}
\]
Substitution and division by $m$ give
\[
 \varepsilon_n-\varepsilon_{n+1}
 =(m-1)p_n-\frac{m-1}{2}\varepsilon_n^2
   +O(p_n^2+\varepsilon_n^3).
\]
The error is $o(n^{-2})$ because $p_n=O(n^{-2})$ and $\varepsilon_n=O(n^{-1})$. The limits of $n^2p_n$ and $n\varepsilon_n$ now give
\[
 n^2(\varepsilon_n-\varepsilon_{n+1})\longrightarrow
 (m-1)\frac4{(m-1)^2}-\frac{m-1}{2}\frac4{(m-1)^2}
 =\frac2{m-1}.
\]
This is the first limit in \eqref{eq:main2}.
\end{proof}

\section{A pointwise bound}\label{sec:pointwise}
The first step towards Proposition~\ref{prop:smoothing} is a pointwise bound on $\rho_n$.
%
We state this bound and an estimate for $\phi_{s,N}$ used in its proof, then prove them in that order.
%
Throughout this section, set
\[
 K_m=2m(m-1)+1,\qquad
 \phi_{s,N}=\frac{q_{s,N-s+1}}{C_s}\quad(N\ge s),
\]
where $C_s$ is the product defined in \eqref{eq:c}.

\begin{proposition}\label{prop:pointwise}
For every $n,j\ge0$,
\begin{equation}\label{eq:pointwise}
 \rho_n(j)\le C\frac{n+1}{(n+j+1)^3}.
\end{equation}
This is equivalent to \eqref{eq:sup}.
\end{proposition}

To prove Proposition~\ref{prop:pointwise}, we use the following estimate, whose proof is deferred to the end of this section.
\begin{proposition}\label{prop:far}
For integers $s\ge0$ and $N>s+K_m$,
\begin{equation}\label{eq:far}
 \phi_{s,N}\le\frac{C(s+1)}{N^3}.
\end{equation}
\end{proposition}

\begin{proof}[Proof of Proposition~\ref{prop:pointwise}]

For $j>K_m$, we have $n+j>n+K_m$. Thus \eqref{eq:far} and $C_n\le1$ give
\[
 \rho_n(j)=(m-1)C_n\phi_{n,n+j}
 \le C\frac{n+1}{(n+j)^3}
 \le C\frac{n+1}{(n+j+1)^3}.
\]

It remains to treat $0\le j\le K_m$. Set $r_0=K_m+2$ and first assume $n\ge2r_0$.
For sequences on $\Zp$, define
\[
 (Tf)(j)=
 \begin{cases}0,&j=0,\\ f(j-1),&j\ge1,\end{cases}
 \qquad
 (f*g)(j)=\sum_{\ell=0}^j f(\ell)g(j-\ell),
\]
and write $f^{*r}$ for the $r$-fold convolution power. Write $Q_s(z)=H_s(z)-x_s$. Recall from \eqref{eq:tilt} that
\[
\begin{aligned}
 zD_sH_{s+1}(z)
 &=H_s(z)^m+(mz-1)x_s^m\\
 &=mz x_s^m+mx_s^{m-1}Q_s(z)
   +\sum_{r=2}^m\binom mr x_s^{m-r}Q_s(z)^r.
\end{aligned}
\]
Since
\[
 Q_s(z)=\frac{z}{m-1}\sum_{j\ge0}\rho_s(j)z^j,\qquad
 [z^{j+2}]Q_s(z)^r
 =\frac{(T^{r-2}\rho_s^{*r})(j)}{(m-1)^r}\quad(r\ge2),
\]
extracting the coefficient of $z^{j+2}$ in the recalled identity and using $c_s=mx_s^{m-1}/D_s$ yields
\begin{equation}\label{eq:rhor}
 \rho_{s+1}=c_sS\rho_s+\sum_{r=2}^m d_{r,s}T^{r-2}\rho_s^{*r},\qquad
 d_{r,s}=\frac{\binom mr x_s^{m-r}}{D_s(m-1)^{r-1}}.
\end{equation}
Using $C_{s+1}=c_sC_s$, iteration from $n-r_0$ to $n$ gives
\[
\begin{aligned}
 \rho_n(j)
 &=\frac{C_n}{C_{n-r_0}}\rho_{n-r_0}(j+r_0)\\
 &\quad+\sum_{s=n-r_0}^{n-1}\frac{C_n}{C_{s+1}}
       \sum_{r=2}^m d_{r,s}
       \bigl(S^{n-1-s}T^{r-2}\rho_s^{*r}\bigr)(j).
\end{aligned}
\]
Recall from \eqref{eq:c}--\eqref{eq:C} that $0<c_i\le1$ and $C_n=\prod_{i=0}^{n-1}c_i$. Also, $x_s=q_{s,0}\in[0,1]$ and $D_s=1+(m-1)x_s^m\ge1$. Thus
\[
\begin{aligned}
 0<\frac{C_n}{C_s}
 &=\prod_{i=s}^{n-1}c_i\le1,
 &&n-r_0\le s\le n,\\
 0\le d_{r,s}
 &\le\frac{\binom mr}{(m-1)^{r-1}}\le2^m,
 &&2\le r\le m.
\end{aligned}
\]
The last inequality uses $(m-1)^{r-1}\ge1$ and $\binom mr\le\sum_{k=0}^m\binom mk=2^m$, giving a bound uniform in $r$ and $s$.
 Since $n+j>(n-r_0)+K_m$, applying \eqref{eq:far} with $s=n-r_0$ and $N=n+j$ gives
\[
\begin{aligned}
 \frac{C_n}{C_{n-r_0}}\rho_{n-r_0}(j+r_0)
 &=(m-1)C_n\phi_{n-r_0,n+j}\\
 &\le C\frac{n-r_0+1}{(n+j)^3}
 \le\frac{C}{(n+1)^2}.
\end{aligned}
\]
To bound the terms in the double sum, recall from \eqref{eq:weightedmass} that, with $L_s=s+1$,
\[
 \norm{\rho_s}_{1,3,L_s}
 =\sum_{k\ge0}\left(1+\frac{k}{s+1}\right)^3\rho_s(k)
 \le\frac{C}{s+1}.
\]
The operators $S$ and $T$ are contractions in the supremum norm. Since $s\ge n-r_0\ge n/2$, nonnegativity and this bound give
\[
\begin{aligned}
 \bigl\|S^{n-1-s}T^{r-2}\rho_s^{*r}\bigr\|_\infty
 &\le\norm{\rho_s^{*r}}_\infty
 \le\norm{\rho_s^{*r}}_1
 =\norm{\rho_s}_1^r\\
 &\le\norm{\rho_s}_{1,3,L_s}^r
 \le\frac{C}{(s+1)^r}
 \le\frac{C}{(n+1)^2},
 \qquad 2\le r\le m.
\end{aligned}
\]
The double sum has $r_0(m-1)$ terms. Hence, for $0\le j\le K_m$,
\[
\begin{aligned}
 \rho_n(j)
 &\le\frac{C\{1+r_0(m-1)\}}{(n+1)^2}
 \le\frac{C}{(n+1)^2}\\
 &=C\left(\frac{n+j+1}{n+1}\right)^3
       \frac{n+1}{(n+j+1)^3}
 \le C\frac{n+1}{(n+j+1)^3},
\end{aligned}
\]
where $(n+j+1)/(n+1)\le1+K_m$.

Finally, for $0\le n<2r_0$ and every $j\ge0$, \eqref{eq:weightedmass} directly gives
\[
 \left(1+\frac{j}{n+1}\right)^3\rho_n(j)
 \le\norm{\rho_n}_{1,3,L_n}
 \le\frac{C}{n+1}
 \le\frac{2Cr_0}{(n+1)^2}.
\]
Rearranging proves \eqref{eq:pointwise} also for these generations.
\end{proof}

\subsection{A change of variables in the recursion}
We derive a recursion for $F_s$ and bounds on the moments and tails of its coefficients, which will be used to estimate $A_s$.

Recall $Q_s(z)=H_s(z)-x_s$ and define
\[
 F_s(z)=\sum_{N\ge s}\phi_{s,N}z^N.
\]
For $k\ge1$, the linear term in the recursion for $q_{s+1,k}$ is $c_sq_{s,k+1}$, and both probabilities correspond to $N=s+k$. Since $C_{s+1}=c_sC_s$, division by $C_{s+1}$ cancels the coefficient $c_s$. Put
\[
 b_s=\phi_{s,s}=\frac{q_{s,1}}{C_s},\qquad
 \kappa_s=\frac{C_s}{x_s},\qquad
 \alpha_{r,s}=\frac1m\binom mr\kappa_s^{r-1},
\]
and define
\begin{equation}\label{eq:arrivalA}
 A_s(z)=\sum_{r=2}^m\alpha_{r,s}z^{(r-2)(1-s)}F_s(z)^r.
\end{equation}
Although the factor $z^{(r-2)(1-s)}$ may involve a negative power of $z$, each summand has minimum degree $2s+r-2$, hence is an ordinary power series with nonnegative coefficients.

\begin{lemma}\label{lem:arrival}
The exact recursions are
\begin{align}
 F_{s+1}(z)&=F_s(z)-b_sz^s+z^{1-s}A_s(z),\label{eq:arrivalrenewal}\\
 \kappa_{s+1}&=\frac{\kappa_s}{1+\kappa_sb_s}.\label{eq:kappa}
\end{align}
In particular, for $N\ge s$,
\begin{equation}\label{eq:arrivalDuhamel}
 \phi_{s,N}=\phi_{0,N}+\sum_{j=0}^{s-1}[z^{N+j-1}]A_j(z).
\end{equation}
\end{lemma}
\begin{proof}
Multiplying \eqref{eq:tilt} by $zD_s$ and expanding $H_s(z)=x_s+Q_s(z)$ gives
\[
\begin{aligned}
 zD_sH_{s+1}(z)
 &= (x_s+Q_s(z))^m+(mz-1)x_s^m\\
 &= mz x_s^m+mx_s^{m-1}Q_s(z)
 +\sum_{r=2}^m\binom mr x_s^{m-r}Q_s(z)^r.
\end{aligned}
\]
For $k\ge1$, the coefficient of $z^{k+1}$ on the left is $D_sq_{s+1,k}$. On the right, $mz x_s^m$ contributes nothing because $k+1\ge2$, while $[z^{k+1}]Q_s(z)=q_{s,k+1}$. Dividing by $D_s$ and using $c_s=mx_s^{m-1}/D_s$, we obtain
\[
 q_{s+1,k}=c_sq_{s,k+1}+
 \frac1{D_s}[z^{k+1}]\sum_{r=2}^m\binom mr x_s^{m-r}Q_s(z)^r.
\]
By definition, $F_{s+1}(z)=C_{s+1}^{-1}\sum_{k\ge1}q_{s+1,k}z^{s+k}$. We therefore multiply the preceding recursion by $z^{s+k}/C_{s+1}$ and sum over $k\ge1$. Since $C_{s+1}=c_sC_s$, the linear term becomes
\[
 \frac{c_s}{C_{s+1}}\sum_{k\ge1}q_{s,k+1}z^{s+k}
 =\frac1{C_s}\sum_{\ell\ge2}q_{s,\ell}z^{s+\ell-1}
 =F_s(z)-b_sz^s.
\]
The subtracted term is precisely the $\ell=1$ term of $F_s$, which has degree $s$.

For each $r\ge2$, the series $Q_s(z)^r$ has no terms of degree zero or one. Hence the sum over $k\ge1$ includes all its coefficients. Substituting $Q_s(z)=C_sz^{1-s}F_s(z)$, the order-$r$ contribution is
\[
\begin{aligned}
 \frac{\binom mr x_s^{m-r}}{D_sC_{s+1}}
   \sum_{k\ge1}z^{s+k}[z^{k+1}]Q_s(z)^r
 &=\frac{\binom mr x_s^{m-r}C_s^r}{D_sC_{s+1}}
   z^{s-1+r(1-s)}F_s(z)^r\\
 &=\frac1m\binom mr\left(\frac{C_s}{x_s}\right)^{r-1}
   z^{(r-1)(1-s)}F_s(z)^r\\
 &=z^{1-s}\alpha_{r,s}z^{(r-2)(1-s)}F_s(z)^r,
\end{aligned}
\]
where we used $D_sC_{s+1}=mx_s^{m-1}C_s$. Summing over $r=2,\ldots,m$ gives $z^{1-s}A_s(z)$ and proves \eqref{eq:arrivalrenewal}.

For \eqref{eq:kappa}, use $x_{s+1}=c_s(x_s+q_{s,1})$ together with $C_{s+1}=c_sC_s$. Since $q_{s,1}/x_s=\kappa_sb_s$, this gives
\[
 \kappa_{s+1}=\frac{C_{s+1}}{x_{s+1}}=\frac{C_s}{x_s+q_{s,1}}
 =\frac{\kappa_s}{1+\kappa_sb_s}.
\]
Finally, fix $N\ge s$. For every $0\le j<s$, the boundary term $b_jz^j$ in \eqref{eq:arrivalrenewal} has degree $j<N$, so taking the coefficient of $z^N$ yields
\[
 \phi_{j+1,N}-\phi_{j,N}
 =[z^N]\bigl(z^{1-j}A_j(z)\bigr)
 =[z^{N+j-1}]A_j(z).
\]
Summing over $j=0,\ldots,s-1$ telescopes the left-hand side to $\phi_{s,N}-\phi_{0,N}$, proving \eqref{eq:arrivalDuhamel}.
\end{proof}

\begin{lemma}[Moment and tail bounds]\label{lem:Fmoment}
For integers $s\ge0$ and $1\le r\le m^2$,
\begin{equation}\label{eq:Fmoment}
 F_s(1)\le\frac C{s+1},\qquad
 \sum_{k\ge0}(k+r)^3[z^k]F_s(z)^r\le C_r(s+1)^{3-r}.
\end{equation}
Consequently, for $K\ge1$,
\begin{equation}\label{eq:Ftail}
 \sum_{k\ge K}[z^k]F_s(z)^r\le C_r(s+1)^{3-r}K^{-3}.
\end{equation}
The constants are independent of $s$ and $K$.
\end{lemma}
\begin{proof}
By the definition of $\phi_{s,N}$, the bound on $\theta_s$ in \eqref{eq:theta}, and $C_s\ge C_\infty>0$ from \eqref{eq:C},
\[
 F_s(1)=\sum_{N\ge s}\phi_{s,N}
 =\frac1{C_s}\sum_{k\ge1}q_{s,k}
 =\frac{\theta_s}{C_s}\le\frac C{s+1}.
\]
The same change of variables $N=s+k-1$, together with $(s+k)^3\le4(s^3+k^3)$, gives
\[
\begin{aligned}
 \sum_{N\ge s}(N+1)^3\phi_{s,N}
 &=\frac1{C_s}\sum_{k\ge1}(s+k)^3q_{s,k}\\
 &\le\frac4{C_s}\left(s^3\theta_s+\sum_{k\ge1}k^3q_{s,k}\right)
 \le C(s+1)^2,
\end{aligned}
\]
where the last inequality uses \eqref{eq:theta} and the third-moment bound in \eqref{eq:lowmoments}.

To estimate $F_s(z)^r$, expand its coefficients as products of the $\phi_{s,N_i}$. After summing over all $k\ge0$, each index $N_1,\ldots,N_r$ ranges independently over $s,s+1,\ldots$. All terms are nonnegative, so the sums can be rearranged. Using
\[
 (N_1+\cdots+N_r+r)^3
 =\left(\sum_{i=1}^r(N_i+1)\right)^3
 \le r^2\sum_{i=1}^r(N_i+1)^3,
\]
we obtain
\[
\begin{aligned}
 \sum_{k\ge0}(k+r)^3[z^k]F_s(z)^r
  &=\sum_{N_1,N_2,\cdots,N_r\ge s}
   (N_1+\cdots+N_r+r)^3\prod_{i=1}^r\phi_{s,N_i}\\
  &\le r^2\sum_{i=1}^r\sum_{N_1,N_2,\cdots,N_r\ge s}
   (N_i+1)^3\prod_{\ell=1}^r\phi_{s,N_\ell}\\
 &=r^3\left(\sum_{N\ge s}(N+1)^3\phi_{s,N}\right)F_s(1)^{r-1}\\
 &\le C_r(s+1)^2(s+1)^{-(r-1)}
 =C_r(s+1)^{3-r}.
\end{aligned}
\]
Indeed, for each fixed $i$, summing over the other $r-1$ indices gives $F_s(1)^{r-1}$, and the $r$ choices of $i$ give identical contributions. This proves \eqref{eq:Fmoment}.

Finally, $k\ge K$ implies $(k+r)^3\ge K^3$. Nonnegativity of the coefficients and \eqref{eq:Fmoment} therefore give
\[
 \sum_{k\ge K}[z^k]F_s(z)^r
 \le K^{-3}\sum_{k\ge K}(k+r)^3[z^k]F_s(z)^r
 \le C_r(s+1)^{3-r}K^{-3},
\]
which is \eqref{eq:Ftail}.
\end{proof}

\subsection{Differences between successive generations}
We bound the coefficients of $A_{s+1}-A_s$ and sum the resulting estimates with $N$ fixed to prove Proposition~\ref{prop:far}.

Recall that $K_m=2m(m-1)+1$.
\begin{lemma}\label{lem:sourcedifference}
For integers $N>s+K_m$,
\begin{equation}\label{eq:sourcedifference}
 \sum_{k\ge N+s}\bigl|[z^k](A_{s+1}-A_s)\bigr|
 \le \frac C{N^3}\{1+b_s(s+1)^2\}.
\end{equation}
\end{lemma}
\begin{proof}
We first control the coefficients $\alpha_{r,s}$. The bounds in \eqref{eq:C} and $\inf_s x_s>0$ imply that $\kappa_s=C_s/x_s$ and all $\alpha_{r,s}$ are bounded. By \eqref{eq:kappa},
\[
 \alpha_{r,s+1}=\alpha_{r,s}(1+\kappa_sb_s)^{-(r-1)}.
\]
Since $0\le1-(1+t)^{-(r-1)}\le(r-1)t$ for $t\ge0$, it follows that
\[
 0\le\alpha_{r,s}-\alpha_{r,s+1}\le Cb_s
 \qquad(2\le r\le m).
\]

In the definition of $A_{s+1}-A_s$, change the coefficient, the power of $z$, and the power of $F_s$ one at a time. Adding and subtracting the corresponding intermediate terms gives
\begin{align}\label{eq:threechanges}
 A_{s+1}-A_s=\sum_{r=2}^m\bigl[&
  (\alpha_{r,s+1}-\alpha_{r,s})z^{(r-2)(1-s)}F_s^r\notag\\
 &+\alpha_{r,s+1}\{z^{-(r-2)s}-z^{(r-2)(1-s)}\}F_s^r\notag\\
 &+\alpha_{r,s+1}z^{-(r-2)s}(F_{s+1}^r-F_s^r)\bigr].
\end{align}
We estimate the coefficient tail of each line by the triangle inequality.

\emph{Coefficient and shift changes.}
Multiplication by $z^a$ changes a tail starting at $N+s$ into one starting at $N+s-a$. For the two powers of $z$ in the first two lines of \eqref{eq:threechanges}, these starting indices are
\[
\begin{aligned}
 N+s-(r-2)(1-s)&=N+(r-1)s-(r-2)>N/2,\\
 N+s+(r-2)s&=N+(r-1)s\ge N.
\end{aligned}
\]
Here we used $r\le m$ and $N>s+K_m$. Thus \eqref{eq:Ftail} and the coefficient bound above show that, for each $r$, the first line contributes at most
\[
 \frac{Cb_s(s+1)^{3-r}}{N^3}
 \le\frac{Cb_s(s+1)^2}{N^3}.
\]
In the second line, the two powers of $z$ agree when $r=2$, so that term vanishes. For $r\ge3$, bound their tails separately using \eqref{eq:Ftail}; the resulting bound is $C(s+1)^{3-r}N^{-3}\le CN^{-3}$.

\emph{The boundary term in the power difference.}
Since $b_sz^s$ is exactly the degree-$s$ term of $F_s$, the series $F_s-b_sz^s$ has nonnegative coefficients, each at most the corresponding coefficient of $F_s$. The recursion \eqref{eq:arrivalrenewal} and the binomial formula give
\[
\begin{aligned}
 F_{s+1}^r-F_s^r
 ={}&(F_s-b_sz^s)^r-F_s^r\\
 &+\sum_{j=1}^r\binom rj(F_s-b_sz^s)^{r-j}
       (z^{1-s}A_s)^j.
\end{aligned}
\]
The first difference has nonpositive coefficients. To bound their absolute values, use
\[
 F_s^r-(F_s-b_sz^s)^r
 =b_sz^s\sum_{j=0}^{r-1}F_s^{r-1-j}(F_s-b_sz^s)^j.
\]
Each of the $r$ products on the right has coefficients bounded by those of $F_s^{r-1}$. After including the factor $\alpha_{r,s+1}z^{-(r-2)s}$ from \eqref{eq:threechanges}, the boundary contribution to the required tail is therefore at most
\[
\begin{aligned}
 Cb_s\sum_{k\ge N+(r-2)s}[z^k]F_s^{r-1}
 &\le\frac{Cb_s(s+1)^{4-r}}{N^3}\\
 &\le\frac{Cb_s(s+1)^2}{N^3},
\end{aligned}
\]
by \eqref{eq:Ftail} and $r\ge2$.

\emph{The terms containing $A_s$.}
All terms in the sum over $j\ge1$ have nonnegative coefficients. Replacing $(F_s-b_sz^s)^{r-j}$ by $F_s^{r-j}$ can only increase them. Fix $j\in\{1,\ldots,r\}$ and choose orders $\ell_1,\ldots,\ell_j\in\{2,\ldots,m\}$ from the $j$ copies of
\[
 z^{1-s}A_s=\sum_{\ell=2}^m\alpha_{\ell,s}
                  z^{(\ell-1)(1-s)}F_s^\ell.
\]
The total number of factors $F_s$ in this product is
\[
 d=r-j+\sum_{i=1}^j\ell_i
   =r+\sum_{i=1}^j(\ell_i-1).
\]
In particular, $3\le r+j\le d\le r+j(m-1)\le m^2$ and $0\le d-r\le m(m-1)$. The scalar coefficient is bounded by a constant depending only on $m$ and the initial law. Including the outer power $z^{-(r-2)s}$, the remaining series is
\begin{equation*}
 z^{-(r-2)s+\sum_{i=1}^j(\ell_i-1)(1-s)}F_s^d
 =z^{d-r+(2-d)s}F_s^d.
\end{equation*}
Consequently its tail starting at $N+s$ corresponds to the tail of $F_s^d$ starting at
\[
 N+s-\{d-r+(2-d)s\}
 =N+(d-1)s-(d-r)
 \ge N-m(m-1)>N/2.
\]
Since $3\le d\le m^2$, \eqref{eq:Ftail} now gives
\begin{equation*}
\begin{aligned}
 \sum_{k\ge N+s}[z^k]\bigl(z^{d-r+(2-d)s}F_s^d\bigr)
 &=\sum_{k\ge N+(d-1)s-(d-r)}[z^k]F_s^d\\
 &\le\frac{C(s+1)^{3-d}}{N^3}\le\frac C{N^3}.
\end{aligned}
\end{equation*}
There are only finitely many choices of $r$, $j$, and $\ell_1,\ldots,\ell_j$, with their number depending only on $m$. Combining this estimate with the bounds for the other terms in \eqref{eq:threechanges} proves \eqref{eq:sourcedifference}.
\end{proof}

We now sum the estimate for $A_{v+1}-A_v$ with $N$ fixed to establish \eqref{eq:far}.
\begin{proof}[Proof of Proposition~\ref{prop:far}]
Since $C_0=1$, the initial term is controlled directly by the third moment:
\[
 \phi_{0,N}=q_{0,N+1}
 \le\frac1{(N+1)^3}\sum_{k\ge1}k^3q_{0,k}
 \le\frac C{N^3}.
\]
Also, \eqref{eq:arrivalA} at $s=0$ and \eqref{eq:Ftail} give
\[
\begin{aligned}
 \sum_{k\ge N-1}[z^k]A_0(z)
 &=\sum_{r=2}^m\alpha_{r,0}
       \sum_{k\ge N-r+1}[z^k]F_0(z)^r\\
 &\le C\sum_{r=2}^m(N-r+1)^{-3}
 \le\frac C{N^3},
\end{aligned}
\]
where $N-r+1>N/2$ follows from $N>K_m$ and $r\le m$.

Now let $s\ge1$. Substitute $A_j=A_0+\sum_{v=0}^{j-1}(A_{v+1}-A_v)$ into \eqref{eq:arrivalDuhamel} and interchange the finite sums:
\[
\begin{aligned}
 \phi_{s,N}
 &=\phi_{0,N}+\sum_{j=0}^{s-1}[z^{N+j-1}]A_0
 +\sum_{j=0}^{s-1}\sum_{v=0}^{j-1}
          [z^{N+j-1}](A_{v+1}-A_v)\\
 &=\phi_{0,N}+\sum_{k=N-1}^{N+s-2}[z^k]A_0
+\sum_{v=0}^{s-2}\sum_{k=N+v}^{N+s-2}
          [z^k](A_{v+1}-A_v).
\end{aligned}
\]
The last equality uses $v+1\le j\le s-1$ and $k=N+j-1$. Empty sums are zero. For $0\le v\le s-2$, we have $N>v+K_m$, so \eqref{eq:sourcedifference} yields
\[
\begin{aligned}
 \left|\sum_{k=N+v}^{N+s-2}[z^k](A_{v+1}-A_v)\right|
 &\le\sum_{k\ge N+v}\bigl|[z^k](A_{v+1}-A_v)\bigr|\\
 &\le\frac C{N^3}\{1+b_v(v+1)^2\}.
\end{aligned}
\]
Combining these estimates gives
\[
\begin{aligned}
 \phi_{s,N}
 &\le\frac C{N^3}
       +\frac C{N^3}\sum_{v=0}^{s-2}\{1+b_v(v+1)^2\}\\
 &\le\frac C{N^3}\left\{s+\sum_{v=0}^{s-1}b_v(v+1)^2\right\}.
\end{aligned}
\]
Finally, $b_v=q_{v,1}/C_v\le C_\infty^{-1}q_{v,1}$ and $q_{v,1}=T_v-T_{v+1}$. Summation by parts and \eqref{eq:atomtail} therefore give
\[
\begin{aligned}
 \sum_{v=0}^{s-1}(v+1)^2b_v
 &\le C_\infty^{-1}\sum_{v=0}^{s-1}(v+1)^2(T_v-T_{v+1})\\
 &=C_\infty^{-1}\left\{T_0+\sum_{v=1}^{s-1}(2v+1)T_v-s^2T_s\right\}\\
 &\le C\left\{1+\sum_{v=1}^{s-1}\frac{2v+1}{v+1}\right\}
 \le C(s+1).
\end{aligned}
\]
This proves \eqref{eq:far} for $s\ge1$. The case $s=0$ is the initial-term estimate above.
\end{proof}

\section{Proof of Proposition~\ref{prop:smoothing}}\label{sec:smoothing}
Proposition~\ref{prop:pointwise} gives the required pointwise bound.
%
To control spatial differences, we compare the recursions for $\rho_s$ and $\rho_s*\rho_s$.
%
We sum the resulting differences from a time $a$ that will be chosen according to the spatial increment $h$. This will give the bound in \eqref{eq:spmod} for every $1\le h\le n+1$.
\subsection{Preliminaries}


To prove \eqref{eq:spmod}, we iterate \eqref{eq:rhor}
from time $a$ to time $n$. 
% The contribution of $\rho_a$ can be bounded
% using Proposition~\ref{prop:pointwise}. We will estimate the spatial
% difference of the remaining contribution by comparing the quadratic
% convolutions at successive times.

Recall the coefficients $d_{r,s}$ from \eqref{eq:rhor}, and put
\[
 d_s=d_{2,s},\qquad B_s=\rho_s*\rho_s,\qquad
 E_s=\sum_{r=3}^m d_{r,s}T^{r-2}\rho_s^{*r},
\]
where $E_s=0$ if $m=2$. Then \eqref{eq:rhor} becomes
\begin{equation}\label{eq:rhoshort}
 \rho_{s+1}=c_sS\rho_s+d_sB_s+E_s.
\end{equation}
For integers $n\ge4$ and $0\le a\le n-2$, define
\begin{align*}
 I_{a,n}&=\frac{C_n}{C_a}S^{n-a}\rho_a,\\
 V_{a,n}&=\sum_{s=a}^{n-1}\frac{C_n}{C_{s+1}}
                   S^{n-1-s}(d_sB_s+E_s).
\end{align*}
Since $C_{s+1}=c_sC_s$, equation~\eqref{eq:rhoshort} gives
\[
 \frac{S^{n-1-s}\rho_{s+1}}{C_{s+1}}
 -\frac{S^{n-s}\rho_s}{C_s}
 =\frac{S^{n-1-s}(d_sB_s+E_s)}{C_{s+1}}.
\]
Summing over $a\le s\le n-1$ and multiplying by $C_n$, we obtain
\begin{equation}\label{eq:split}
 \rho_n=I_{a,n}+V_{a,n}.
\end{equation}
We first bound $I_{a,n}$. By \eqref{eq:C} and \eqref{eq:pointwise}, 
\[
 \begin{aligned}
 I_{a,n}(j)
 &=\frac{C_n}{C_a}\rho_a(j+n-a)
 \le C\frac{a+1}{(n+j+1)^3},\\
 \left(1+\frac j{L_n}\right)^3 I_{a,n}(j)
 &\le \frac{C(a+1)}{L_n^3}.
 \end{aligned}
\]
Taking the supremum over $j\ge0$ gives
\begin{equation}\label{eq:earlyI}
 \norm{I_{a,n}}_{\infty,3,L_n}
 \le\frac{C(a+1)}{L_n^3}.
\end{equation}
For $L\ge1$ and an integer $k\ge0$, the inequality
$(1+j/L)^3\le(1+(j+k)/L)^3$ implies
\[
 \norm{S^kf}_{\infty,3,L}\le\norm f_{\infty,3,L}.
\]
Consequently, for every integer $h\ge1$,
\[
 \begin{aligned}
 \norm{I_{a,n}-S^hI_{a,n}}_{\infty,3,L_n}
 &\le2\norm{I_{a,n}}_{\infty,3,L_n},\\
 \norm{V_{a,n}-S^hV_{a,n}}_{\infty,3,L_n}
 &=\left\|\sum_{\ell=0}^{h-1}S^\ell\Delta V_{a,n}
       \right\|_{\infty,3,L_n}
 \le h\norm{\Delta V_{a,n}}_{\infty,3,L_n}.
 \end{aligned}
\]
Combining these inequalities with \eqref{eq:split} and \eqref{eq:earlyI},
we reduce the spatial estimate to a bound on $\Delta V_{a,n}$:
\begin{equation}\label{eq:spatialreduction}
 \norm{\rho_n-S^h\rho_n}_{\infty,3,L_n}
 \le\frac{C(a+1)}{L_n^3}
       +h\norm{\Delta V_{a,n}}_{\infty,3,L_n}.
\end{equation}
It remains to prove \eqref{eq:lateV}, namely
\[
 \norm{\Delta V_{a,n}}_{\infty,3,L_n}
 \le \frac{C}{L_n^3}
       \left(1+\log\frac{L_n}{a+1}\right).
\]
For $1\le h\le n/2$, this estimate and
\eqref{eq:spatialreduction} give \eqref{eq:spmod} with $a=h$.
For $n/2<h\le L_n$, \eqref{eq:spmod} follows directly
from \eqref{eq:pointwise}.

To prepare the estimate of $\Delta V_{a,n}$, we now bound $B_s$ and $E_s$
and derive a recursion relating $B_{s+1}$ to $B_s$.
The required bounds follow from two weighted convolution inequalities.
Fix $L\ge1$.
For every $i,j\in\Zp$,
\begin{equation}\label{eqn:weight}
 \left(1+\frac{i+j}{L}\right)^3
 \le
 \left(1+\frac iL\right)^3
 \left(1+\frac jL\right)^3.
\end{equation}
Let $f,g$ be sequences on $\Zp$ with finite weighted sum norms.
The triangle inequality and \eqref{eqn:weight} give
\[
 \begin{aligned}
 \norm{f*g}_{1,3,L}
 &\le
 \sum_{k\ge0}\sum_{i=0}^{k}
 \left(1+\frac kL\right)^3
 |f(i)|\,|g(k-i)|\\
 &\le
 \sum_{i,j\ge0}
 \left(1+\frac iL\right)^3|f(i)|
 \left(1+\frac jL\right)^3|g(j)|\\
 &=\norm f_{1,3,L}\norm g_{1,3,L}.
 \end{aligned}
\]
Thus we have
\begin{equation}\label{ineq:con_1}
 \norm{f*g}_{1,3,L}
 \le \norm f_{1,3,L}\norm g_{1,3,L}.
\end{equation}
For the supremum norm, fix $k\ge0$ and use \eqref{eqn:weight} to obtain
\[
 \begin{aligned}
 \left(1+\frac kL\right)^3|(f*g)(k)|
 &\le
 \sum_{i=0}^{k}
 \left(1+\frac iL\right)^3|f(i)|
 \left(1+\frac{k-i}{L}\right)^3|g(k-i)|\\
 &\le
 \norm f_{\infty,3,L}
 \sum_{j=0}^{k}\left(1+\frac jL\right)^3|g(j)|\\
 &\le
 \norm f_{\infty,3,L}\norm g_{1,3,L}.
 \end{aligned}
\]
Taking the supremum over $k$ therefore gives
\begin{equation}\label{ineq:con_inf}
 \norm{f*g}_{\infty,3,L}
 \le \norm f_{\infty,3,L}\norm g_{1,3,L}.
\end{equation}

Now fix $s\ge0$ and recall that $L_s=s+1$.
Equation~\eqref{eq:weightedmass} gives
\[
 \norm{\rho_s}_{1,3,L_s}\le CL_s^{-1}.
\]
Moreover, \eqref{eq:pointwise} implies that, for every $j\ge0$,
\[
 \left(1+\frac j{L_s}\right)^3\rho_s(j)
 \le
 \frac{(L_s+j)^3}{L_s^3}
 \frac{CL_s}{(L_s+j)^3}
 =CL_s^{-2}.
\]
Hence $\norm{\rho_s}_{\infty,3,L_s}\le CL_s^{-2}$.
Starting from these bounds for $r=1$ and applying
\eqref{ineq:con_1} and \eqref{ineq:con_inf} inductively with $L=L_s$,
we obtain, for every fixed positive integer $r$,
\begin{equation}\label{eq:convnorms}
 \begin{aligned}
 \norm{\rho_s^{*r}}_{1,3,L_s}
 &\le \norm{\rho_s}_{1,3,L_s}^{\,r}
 \le C_rL_s^{-r},\\
 \norm{\rho_s^{*r}}_{\infty,3,L_s}
 &\le
 \norm{\rho_s}_{\infty,3,L_s}
 \norm{\rho_s}_{1,3,L_s}^{\,r-1}
 \le C_rL_s^{-r-1}.
 \end{aligned}
\end{equation}
Here $C_r$ is independent of $s$.

In particular, taking $r=2$ in \eqref{eq:convnorms} gives
\[
 \norm{B_s}_{1,3,L_s}\le CL_s^{-2},\qquad
 \norm{B_s}_{\infty,3,L_s}\le CL_s^{-3}.
\]
To estimate $E_s$, we also account for the right shifts in its definition.
For $k,j\ge0$,
\[
 \left(1+\frac{j+k}{L}\right)^3
 \le\left(1+\frac kL\right)^3\left(1+\frac jL\right)^3.
\]
By the definition of $T$, this gives, for $p\in\{1,\infty\}$,
\[
 \norm{T^kf}_{p,3,L}
 \le\left(1+\frac kL\right)^3\norm f_{p,3,L}
 \le(1+k)^3\norm f_{p,3,L}.
\]
The coefficients $d_{r,s}$ are uniformly bounded for $2\le r\le m$
by \eqref{eq:rhor}. Hence \eqref{eq:convnorms} yields
\[
 \begin{aligned}
 \norm{E_s}_{\infty,3,L_s}
 &\le\sum_{r=3}^m d_{r,s}(r-1)^3
                   \norm{\rho_s^{*r}}_{\infty,3,L_s}\\
 &\le C\sum_{r=3}^m L_s^{-r-1}
 \le CL_s^{-4},\\
 \norm{E_s}_{1,3,L_s}
 &\le\sum_{r=3}^m d_{r,s}(r-1)^3
                   \norm{\rho_s^{*r}}_{1,3,L_s}\\
 &\le C\sum_{r=3}^m L_s^{-r}
 \le CL_s^{-3}.
 \end{aligned}
\]
These bounds also hold when $m=2$, since then $E_s=0$.
The next subsection will use the bound on $E_s$ to control its sum in
$V_{a,n}$. For the quadratic terms, we will additionally use their
relation at successive times.

The linear term in \eqref{eq:rhoshort} shifts $\rho_s$ by one site.
Convolving two such terms gives a shift of $B_s$ by two sites, with a
correction at the boundary of $\Zp$. In the definition of $V_{a,n}$,
the exponent $n-1-s$ decreases by one when $s$ increases by one.
The two-site shift in the recursion for $B_s$ will therefore leave
one spatial shift between the quadratic terms propagated to time $n$.
This is what allows their spatial difference to be summed over $s$.
To obtain the exact recursion, put $\beta_s=\rho_s(0)$.
For every $j\ge0$,
\[
 \begin{aligned}
 ((S\rho_s)*(S\rho_s))(j)
 &=\sum_{i=0}^j\rho_s(i+1)\rho_s(j-i+1)\\
 &=B_s(j+2)-2\beta_s\rho_s(j+2),\\
 ((S\rho_s)*B_s)(j)
 &=\sum_{i=0}^j\rho_s(i+1)B_s(j-i)\\
 &=\rho_s^{*3}(j+1)-\beta_sB_s(j+1).
 \end{aligned}
\]
Substituting \eqref{eq:rhoshort} into
$B_{s+1}=\rho_{s+1}*\rho_{s+1}$ first gives
\[
 \begin{aligned}
 B_{s+1}={}&c_s^2(S\rho_s)*(S\rho_s)
       +2c_sd_s(S\rho_s)*B_s+d_s^2\rho_s^{*4}\\
       &+2(c_sS\rho_s+d_sB_s)*E_s+E_s*E_s.
 \end{aligned}
\]
Using the two identities above, we obtain
\begin{equation}\label{eq:Bevolution}
 B_{s+1}=c_s^2S^2B_s+R_s,
\end{equation}
where
\begin{align*}
 R_s={}&-2c_s^2\beta_sS^2\rho_s
       +2c_sd_sS(\rho_s^{*3})
       -2c_sd_s\beta_sSB_s+d_s^2\rho_s^{*4}\\
       &+2(c_sS\rho_s+d_sB_s)*E_s+E_s*E_s.
\end{align*}
We now bound the error $R_s$ in \eqref{eq:Bevolution}.

By \eqref{eq:pointwise}, $\beta_s\le CL_s^{-2}$.
Since $S$ contracts the weighted supremum norm,
\eqref{eq:convnorms} gives
\[
 \begin{aligned}
 \norm{\beta_sS^2\rho_s}_{\infty,3,L_s}
 &\le\beta_s\norm{\rho_s}_{\infty,3,L_s}
 \le CL_s^{-4},\\
 \norm{S(\rho_s^{*3})}_{\infty,3,L_s}
 &\le\norm{\rho_s^{*3}}_{\infty,3,L_s}
 \le CL_s^{-4},\\
 \norm{\beta_sSB_s}_{\infty,3,L_s}
 +\norm{\rho_s^{*4}}_{\infty,3,L_s}
 &\le CL_s^{-5}.
 \end{aligned}
\]
For the terms containing $E_s$, apply \eqref{ineq:con_inf} and the bounds
already proved for $\rho_s$, $B_s$, and $E_s$:
\[
 \begin{aligned}
 \norm{(c_sS\rho_s+d_sB_s)*E_s}_{\infty,3,L_s}
 &\le
 \bigl(c_s\norm{\rho_s}_{\infty,3,L_s}
       +d_s\norm{B_s}_{\infty,3,L_s}\bigr)
       \norm{E_s}_{1,3,L_s}\\
 &\le C(L_s^{-2}+L_s^{-3})L_s^{-3}
 \le CL_s^{-5},\\
 \norm{E_s*E_s}_{\infty,3,L_s}
 &\le\norm{E_s}_{\infty,3,L_s}\norm{E_s}_{1,3,L_s}
 \le CL_s^{-7}.
 \end{aligned}
\]
Since $c_s$ and $d_s$ are bounded and $L_s\ge1$, these estimates imply
\begin{equation}\label{eq:REbounds}
 \norm{R_s}_{\infty,3,L_s}\le CL_s^{-4},\qquad
 \norm{E_s}_{\infty,3,L_s}\le CL_s^{-4},\qquad
 \norm{E_s}_{1,3,L_s}\le CL_s^{-3}.
\end{equation}
We next apply \eqref{eq:Bevolution} to the quadratic terms in $V_{a,n}$
and use \eqref{eq:REbounds} to prove \eqref{eq:lateV}.

\subsection{A bound on the spatial differences}

We now prove \eqref{eq:lateV} for the sequence $V_{a,n}$ defined above.
Fix integers $n\ge4$ and $0\le a\le n-2$.
For $a\le s\le n-1$, put
\[
 P_s=\frac{C_n}{C_{s+1}}=\prod_{i=s+1}^{n-1}c_i,\qquad
 w_s=P_sd_s,\qquad K_s=S^{n-1-s}B_s.
\]
Separate the quadratic and higher-order terms in $V_{a,n}$ by writing
\[
 V_{a,n}=Q_{a,n}+W_{a,n},\qquad
 Q_{a,n}=\sum_{s=a}^{n-1}w_sK_s,\qquad
 W_{a,n}=\sum_{s=a}^{n-1}P_sS^{n-1-s}E_s.
\]
All three sequences $I_{a,n}$, $Q_{a,n}$, and $W_{a,n}$ are nonnegative.

To use the bounds at time $s$ in the weighted norm at time $n$, we need
an estimate for the shifts $S^{n-1-s}$ and $S^{n-2-s}$.
For $0\le s\le n-1$ and an integer $d\ge0$ satisfying $d\ge n-2-s$,
\begin{equation}\label{eq:shiftgain}
 \norm{S^df}_{\infty,3,L_n}
 \le27\left(\frac{L_s}{L_n}\right)^3
       \norm f_{\infty,3,L_s}.
\end{equation}
Indeed, $s+1+d\ge n-1$, and hence
\[
 \frac{(1+j/L_n)^3}{(1+(j+d)/L_s)^3}
 =\left(\frac{L_s}{L_n}\right)^3
  \left(\frac{n+1+j}{s+1+j+d}\right)^3
 \le27\left(\frac{L_s}{L_n}\right)^3.
\]
Multiplying by $(1+(j+d)/L_s)^3|f(j+d)|$ and taking the supremum
proves \eqref{eq:shiftgain}. For $a\le s\le n-2$, \eqref{eq:Bevolution} gives
\[
 K_{s+1}=c_s^2SK_s+S^{n-2-s}R_s.
\]
We rewrite this identity as
\[
 \Delta K_s
 =K_s-c_s^{-2}K_{s+1}
  +c_s^{-2}S^{n-2-s}R_s.
\]
We multiply by $w_s$ and sum over $a\le s\le n-2$.
Using $K_{n-1}=B_{n-1}$, we obtain
\begin{align*}
 \Delta Q_{a,n}
 &=\sum_{s=a}^{n-2}w_s\Delta K_s
   +w_{n-1}\Delta B_{n-1}\\
 &=w_aK_a-w_{n-2}c_{n-2}^{-2}B_{n-1}
   +w_{n-1}\Delta B_{n-1}\\
 &\quad
   +\sum_{s=a+1}^{n-2}
       \bigl(w_s-w_{s-1}c_{s-1}^{-2}\bigr)K_s\\
 &\quad
   +\sum_{s=a}^{n-2}
       w_sc_s^{-2}S^{n-2-s}R_s.
\end{align*}
The differences of the coefficients $w_s$ are controlled as follows.
By \eqref{eq:rhor},
\[
 d_s=\frac{m x_s^{m-2}}{2\{1+(m-1)x_s^m\}}.
\]
The function of $x_s$ on the right has bounded derivative on $[0,1]$.
The recursion for $x_s=q_{s,0}$ in Section~\ref{subsec:transport},
\eqref{eq:C}, and the bound
$q_{s,1}=\rho_s(0)/(m-1)\le CL_s^{-2}$ from
\eqref{eq:pointwise} give
\[
 \begin{aligned}
 |x_{s+1}-x_s|
 &\le(1-c_s)x_s+c_sq_{s,1}
 \le CL_s^{-2},\\
 |d_s-d_{s-1}|
 &\le C|x_s-x_{s-1}|
 \le CL_{s-1}^{-2}
 \le CL_s^{-2}\qquad(s\ge1),
 \end{aligned}
\]
where $L_s/L_{s-1}\le2$ for $s\ge1$.
Also, \eqref{eq:C} gives $|1-c_s|\le CL_s^{-2}$ and
$\inf_s c_s>0$.
Since $P_{s-1}=c_sP_s$,
\[
 |w_s-w_{s-1}c_{s-1}^{-2}|
 =P_s|d_s-c_sc_{s-1}^{-2}d_{s-1}|\le CL_s^{-2}.
\]
We estimate the three types of terms in this telescoping sum separately. By
\eqref{eq:convnorms} and \eqref{eq:shiftgain}, for $a\le s\le n-1$,
\[
 \norm{K_s}_{\infty,3,L_n}
 \le27\left(\frac{L_s}{L_n}\right)^3
       \norm{B_s}_{\infty,3,L_s}
 \le \frac C{L_n^3}.
\]
The coefficients $w_s$ and $c_s^{-1}$ are uniformly bounded. Since
$K_{n-1}=B_{n-1}$ and $S$ contracts the weighted norm, the three boundary
terms therefore satisfy
\[
 \begin{aligned}
 &\norm{w_aK_a}_{\infty,3,L_n}
 +\norm{w_{n-2}c_{n-2}^{-2}B_{n-1}}_{\infty,3,L_n}\\
 &\qquad+\norm{w_{n-1}\Delta B_{n-1}}_{\infty,3,L_n}
 \le \frac C{L_n^3}.
 \end{aligned}
\]
For the sum containing $w_s-w_{s-1}c_{s-1}^{-2}$, the preceding bound on $K_s$ gives
\[
 \begin{aligned}
 &\left\|\sum_{s=a+1}^{n-2}
   (w_s-w_{s-1}c_{s-1}^{-2})K_s\right\|_{\infty,3,L_n}\\
 &\qquad\le\frac C{L_n^3}\sum_{s=a+1}^{n-2}L_s^{-2}
 \le\frac C{L_n^3}.
 \end{aligned}
\]
For the remainder sum, \eqref{eq:REbounds} and \eqref{eq:shiftgain} yield
\[
 \begin{aligned}
 &\left\|\sum_{s=a}^{n-2}
    w_sc_s^{-2}S^{n-2-s}R_s\right\|_{\infty,3,L_n}\\
 &\qquad\le C\sum_{s=a}^{n-2}
    \left(\frac{L_s}{L_n}\right)^3 L_s^{-4}
 =\frac C{L_n^3}\sum_{s=a}^{n-2}\frac1{L_s}.
 \end{aligned}
\]
The harmonic sum in the remainder supplies the logarithm. The integral comparison
\[
 \sum_{s=a}^{n-1}\frac1{L_s}
 \le\frac1{a+1}+\int_{a+1}^{n}\frac{\dd t}{t}
 \le1+\log\frac{L_n}{a+1}
\]
then gives
\begin{equation*}
 \norm{\Delta Q_{a,n}}_{\infty,3,L_n}
 \le \frac C{L_n^3}\left(1+\log\frac{L_n}{a+1}\right).
\end{equation*}
For the terms involving convolutions of order at least three, use $0<P_s\le1$ together with
\eqref{eq:REbounds} and \eqref{eq:shiftgain}:
\[
 \begin{aligned}
 \norm{W_{a,n}}_{\infty,3,L_n}
 &\le\sum_{s=a}^{n-1}P_s
       \norm{S^{n-1-s}E_s}_{\infty,3,L_n}\\
 &\le C\sum_{s=a}^{n-1}
       \left(\frac{L_s}{L_n}\right)^3 L_s^{-4}
 =\frac C{L_n^3}\sum_{s=a}^{n-1}\frac1{L_s}\\
 &\le\frac C{L_n^3}\left(1+\log\frac{L_n}{a+1}\right).
 \end{aligned}
\]
Since $V_{a,n}=Q_{a,n}+W_{a,n}$ and
$\norm{\Delta f}_{\infty,3,L}\le2\norm f_{\infty,3,L}$, the triangle
inequality gives
\[
 \norm{\Delta V_{a,n}}_{\infty,3,L_n}
 \le\norm{\Delta Q_{a,n}}_{\infty,3,L_n}
    +2\norm{W_{a,n}}_{\infty,3,L_n}.
\]
Combining the two estimates proves
\begin{equation}\label{eq:lateV}
 \norm{\Delta V_{a,n}}_{\infty,3,L_n}
 \le \frac C{L_n^3}\left(1+\log\frac{L_n}{a+1}\right).
\end{equation}
The logarithm comes from the sum of $L_s^{-1}$ in the remainder. We will use its dependence on $a$ to prove \eqref{eq:spmod} by choosing $a$ according to the spatial increment $h$.

\subsection{Proof of Proposition~\ref{prop:smoothing}}
By Proposition~\ref{prop:pointwise} and $C_n/C_a\le1$, for $0\le a\le n$,
\[
 I_{a,n}(j)=\frac{C_n}{C_a}\rho_a(j+n-a)
 \le C\frac{a+1}{(n+j+1)^3}.
\]
Hence, in particular for $0\le a\le n/2$,
\begin{equation}\label{eq:earlyIspatial}
 \norm{I_{a,n}}_{\infty,3,L_n}\le C\frac{a+1}{L_n^3}.
\end{equation}
For every integer $h\ge1$, $I-S^h=\sum_{j=0}^{h-1}S^j\Delta$, and $S$ contracts the weighted norm. Equations \eqref{eq:split}, \eqref{eq:lateV}, and \eqref{eq:earlyIspatial} therefore yield
\begin{equation*}
 \norm{\rho_n-S^h\rho_n}_{\infty,3,L_n}
 \le \frac C{L_n^3}\left\{a+1+h\left(1+\log\frac{L_n}{a+1}\right)\right\}
\end{equation*}
for $0\le a\le n/2$. For $1\le h\le n/2$, choose $a=h$. For $n/2<h\le L_n$, the crude bound $2\norm{\rho_n}_{\infty,3,L_n}\le CL_n^{-2}$ implies the same conclusion. The finitely many small $n$ are absorbed into the constant. This proves \eqref{eq:spmod}; setting $h=1$ proves \eqref{eq:gradient}.

For the temporal bound, iterate \eqref{eq:rhor} from $n$ to $n+h$, with $1\le h\le L_n$. By \eqref{eq:C},
\[
 \left|1-\prod_{s=n}^{n+h-1}c_s\right|\le C\frac h{L_n^2}.
\]
For $n\le s<n+h$, the supremum norm of $d_sB_s+E_s$ is $O(L_n^{-3})$ by \eqref{eq:convnorms}. Consequently,
\[
 \norm{\rho_{n+h}-S^h\rho_n}_\infty\le C\frac h{L_n^3}.
\]
Combine this with \eqref{eq:spmod}. This proves \eqref{eq:timemod} and hence Proposition~\ref{prop:smoothing}.

\section{Moment bounds and equations for subsequential limits}\label{sec:limits}
We now prepare the proof of Proposition~\ref{prop:identification}.
%
The bounds and continuity estimates of Proposition \ref{prop:smoothing}, proved in Sections \ref{sec:pointwise} and \ref{sec:smoothing}, give local compactness of the rescaled functions $u_N$, but do not control the third-moment tail.
%
The Markov chain constructed below supplies this missing estimate.
%
We then pass the discrete recursion and its moments to the limit.

\subsection{A Markov chain and its tail bound}
Recall that $h_m$ is the nonnegative cubic polynomial in \eqref{eq:cubic}, and $J_n=\E_{q_n}h_m(K)$ grows at most quadratically, where $K$ is a random variable with law $q_n$. The next construction realizes the probability laws proportional to $h_mq_n$ on a common Markov chain. It is defined only where $h_m$ is positive.

\begin{lemma}\label{lem:spine}
There is a Markov chain $K_n^\sharp$ with marginals
\begin{equation} \label{eq:Markov_marginals}
 \Pp(K^{\sharp}_n = k) = q_n^\sharp(k) := \frac{h_m(k)q_{n,k}}{J_n}
\end{equation}
such that
\begin{equation}\label{eq:increment}
 \E[(K_{n+1}^\sharp-K_n^\sharp)_+]\le4\qquad(n\ge0).
\end{equation}
In particular, for $R>0$,
\begin{equation}\label{eq:sharptail}
 \Pp(K^{\sharp}_n > R) \leq \Pp(K^{\sharp}_0 > R/2) + \frac{8n}{R}.
 % q_n^\sharp(K>R)\le q_0^\sharp(K>R/2)+\frac{8n}{R}.
\end{equation}
\end{lemma}

\begin{proof}
We construct a time-inhomogeneous Markov chain with the marginal distributions \eqref{eq:Markov_marginals} using auxiliary random variables. For $k\in\supp q_n^\sharp$ (so $k>0$), let $Y=(Y_1,\ldots,Y_{m-1})$ be a vector of nonnegative integer-valued random variables and define
\begin{align*}
 A_m(k;y)={}&(k)_3+(k)_2\left(3\sum_i y_i+\chi_m\right)\\
 &+k\left(2\sum_{i<j}y_iy_j+\chi_m\sum_i y_i\right).
\end{align*}
All terms are nonnegative. We then define the law of $Y$ by weighting the product law $q_n^{\otimes(m-1)}$ with the density $A_m(k;y)/h_m(k)$, so that
\[ 
 \Pp(Y = (y_1, \dotsc, y_{m - 1}) \mid K^{\sharp}_n = k) = \frac{A_m(k;y)}{h_m(k)} \prod^{m - 1}_{i = 1} q_{n, y_{i}}.
 \]
The restriction $h_m(k)>0$ applies only to the current state $k$; the auxiliary coordinates $y_i$ range over all nonnegative integers, including zero.
Indeed, independence under $q_n^{\otimes(m-1)}$ and \eqref{eq:qmean} give
\[
 \E_{q_n^{\otimes(m-1)}} A_m(k;Y)
 =(k)_3+(3+\chi_m)(k)_2+2\chi_m k=h_m(k),
\]
where $2\binom{m-1}{2}/(m-1)^2=\chi_m$. Thus the density integrates to one. We define the transition probabilities for $k\in\supp q_n^\sharp$ by
\begin{equation} \label{eq:transition_prob_mult}
 \Pp(K^{\sharp}_{n+1} = \ell \mid K^{\sharp}_n = k) = \Pp(k+Y_1+\cdots+Y_{m-1}-1=\ell).
\end{equation}
For an integer $V\ge0$, the polynomial in \eqref{eq:cubic} satisfies
\begin{equation}\label{eq:h_m(V)_in(V)_k}
 h_m((V-1)_+)=(V)_3+\chi_m(V)_2.
\end{equation}
Now we verify that the transition probabilities are consistent with the marginal distributions \eqref{eq:Markov_marginals}. For $\ell = 0$, it is straightforward to check that \eqref{eq:transition_prob_mult} implies $\Pp(K^{\sharp}_{n + 1} = \ell) = 0$, consistent with \eqref{eq:Markov_marginals}. Below we verify the case $\ell\ge1$. Put $A_m^{(i)}(k_1,\ldots,k_m):=A_m(k_i;(k_j)_{j\ne i})$. For $V=k_1+\cdots+k_m$, \eqref{eq:h_m(V)_in(V)_k} gives
\begin{equation} \label{eq:defn_A^(i)_m}
 \sum_{i=1}^m A^{(i)}_m(k_1, \dotsc, k_m) = (V)_3 + \chi_m (V)_2 = h_m((V-1)_+).
\end{equation}
Hence, from \eqref{eq:transition_prob_mult} and \eqref{eq:defn_A^(i)_m},
\begin{equation*}
\begin{split}
    \Pp(K^{\sharp}_{n + 1} = \ell) = {}& \frac{1}{mJ_n} \sum_{\substack{k_1, \dotsc, k_m \geq 0 \\ k_1 + \dotsb + k_m = \ell + 1}} \left( \sum^m_{i = 1} A^{(i)}_m(k_1, \dotsc, k_m) \right) \prod^m_{j = 1} q_{n, k_j} \\
    = {}& \frac{h_m(\ell)}{mJ_n} \sum_{\substack{k_1, \dotsc, k_m \geq 0 \\ k_1 + \dotsb + k_m = \ell + 1}} \prod^m_{j = 1} q_{n, k_j}.
\end{split} 
\end{equation*}
Since $\ell\ge1$, taking the coefficient of $z^\ell$ in \eqref{eq:tilt} gives
\begin{equation*}
    \sum_{\substack{k_1, \dotsc, k_m \geq 0 \\ k_1 + \dotsb + k_m = \ell + 1}} \prod^m_{j = 1} q_{n, k_j} = D_n q_{n + 1, \ell},
\end{equation*}
and \eqref{eq:J} then yields $\Pp(K^{\sharp}_{n + 1} = \ell) = q^{\sharp}_{n + 1}(\ell)$, as required by \eqref{eq:Markov_marginals}. In particular, when $m=2$, the chain remains in $\{2,3,\ldots\}$.

To prove \eqref{eq:increment}, we put $b= \sum^{\infty}_{k = 0} (k)_2 q_{n, k}$, $g= \sum^{\infty}_{k = 0} (k)_3 q_{n, k}$, $x=x_n = q_{n, 0}$, and $\chi=\chi_m = 1 - 1/(m - 1)$. Computing the expected positive increment using the conditional distribution defined above gives
\begin{equation}\label{eq:incformula}
 \E[(K_{n+1}^\sharp-K_n^\sharp)_+]
 =\frac{x^{m-1}(g+\chi b)+3(m-1)b^2+6\chi b+3\chi^2/(m-1)}
 {g+(3+\chi)b+2\chi/(m-1)}.
\end{equation}
Here is a direct calculation. We let $K, X_1, \dotsc, X_{m - 1}$ be independent random variables with law $q_n$, and let $U=\sum^{m - 1}_{i = 1} X_i$ and $V=\sum_{1 \leq i<j \leq m - 1}X_iX_j$. Then
\begin{equation*}
  \E[(K_{n+1}^\sharp-K_n^\sharp)_+] = \frac{1}{J_n} \E[A_m(K; (X_1, \dotsc, X_{m - 1})) (U - 1)_+],
\end{equation*}
and 
\begin{equation*}
    A_m(K; (X_1, \dotsc, X_{m - 1})) = (K)_3 + \chi(K)_2 + 3(K)_2 U + 2KV + \chi KU.
\end{equation*}
Using
\[
 \E(U-1)_+=x^{m-1},\quad \E[U(U-1)]=(m-1)b+\chi, \quad  \E[V(U-1)]=(m-2)b+\chi^2,
\]
and $J_n = g+(3+\chi)b+2\chi/(m-1)$, we obtain \eqref{eq:incformula}.

Cauchy--Schwarz for the size-biased law gives
$b^2\le(g+b)/(m-1)$. Since $0\le\chi<1$ and $x\le1$, the numerator in \eqref{eq:incformula} is at most
\[
 4g+(3+7\chi)b+\frac{3\chi^2}{m-1}
 \le4\left\{g+(3+\chi)b+\frac{2\chi}{m-1}\right\}.
\]
This proves \eqref{eq:increment}, including the case $m=2$, where $\chi=0$.

Pathwise, $K_n^\sharp\le K_0^\sharp+\sum_{s<n}(K_{s+1}^\sharp-K_s^\sharp)_+$. A union bound and Markov's inequality yield \eqref{eq:sharptail}. The estimate only uses the tail of $K_0^\sharp$; its first moment need not be finite.
\end{proof}

Combining Lemma~\ref{lem:spine} with the cubic moment bound gives uniform integrability under the scaling used in \eqref{eq:rescaled}.
\begin{corollary}\label{cor:thirdtail}
For $R\ge2$,
\begin{equation}\label{eq:thirdtail}
 \sum_{k>R}k^3q_{n,k}
 \le C(n+1)^2\left\{ \Pp(K^\sharp_0>R/2)+\frac nR\right\}.
\end{equation}
For every $0<\eta<T<\infty$,
\begin{equation}\label{eq:UI}
 \lim_{A\to\infty}\limsup_{N\to\infty}
 \sup_{\eta N\le n\le TN}N^{-2}\sum_{k>AN}k^3q_{n,k}=0.
\end{equation}
\end{corollary}
\begin{proof}
For integers $k>R\ge2$, we have $k\ge3$ and
\[
 h_m(k)\ge k^3-k\ge\frac89k^3,
 \qquad k^3\le\frac98h_m(k)\qquad(m\ge2).
\]
Multiply \eqref{eq:sharptail} by $(9/8)J_n$ and use \eqref{eq:J}. For \eqref{eq:UI}, divide by $N^2$, take the supremum, first let $N\to\infty$ to remove the tail of the single fixed probability law $q_0^\sharp$, and then let $A\to\infty$. The order of limits is essential.
\end{proof}

\subsection{Proof of Proposition~\ref{prop:limit}} \label{subsec:prop:limit_proof}
Let $u$ be a locally uniform subsequential limit of the interpolated functions in \eqref{eq:rescaled}. Such limits exist by Proposition~\ref{prop:smoothing}, as shown at the start of the proof of Theorem~\ref{thm:profile}. We work along the corresponding subsequence throughout this section.

\begin{proposition}\label{prop:limit}
The limit $u$ has moments
\[
 A_r(t)=\int_0^\infty x^ru(t,x)\dd x\qquad(0\le r\le3)
\]
satisfying
\begin{equation}\label{eq:Amoment}
 A_0(t)\le C/t,\qquad A_1(t)=1,\qquad A_2(t)\le Ct,\qquad A_3(t)\le Ct^2.
\end{equation}
The third moment also satisfies the small-time tail estimate
\begin{equation}\label{eq:limitthirdtail}
 \int_{x>\delta}x^3u(t,x)\dd x\le C\frac{t^3}{\delta}\qquad(t,\delta>0).
\end{equation}
The function $u$ satisfies the following integral equation, which is the mild formulation of $\partial_t u = \partial_x u + \frac{1}{2} u*u$:
\begin{equation}\label{eq:mild}
 u(t,x)=u(t_0,x+t-t_0)+\frac12\int_{t_0}^t(u(s,\cdot)*u(s,\cdot))(x+t-s)\dd s
\end{equation}
for every $0<t_0<t$ and $x\ge0$.
\end{proposition}
\begin{proof}
Let $\{ N_a \}$ be a subsequence of natural numbers such that $u_{N_a}$ converges locally uniformly. For a grid time with $n/N_a \to t>0$, consider the measure
\[
 \mu_{N_a,t}=N_a \sum_{j\ge0}\rho_n(j)\delta_{j/N_a}.
\]
Its local limit is $u(t,x)\dd x$, by locally uniform convergence and the Riemann-sum normalization. Its $r$th moment is $N_a^{1-r}\sum_jj^r\rho_n(j)$. The coarse bounds and Corollary~\ref{cor:thirdtail} give uniform integrability through order three on every compact time interval contained in $(0, +\infty)$. Hence moments pass to the limit:
\begin{equation*}
 A_r(t)=\lim_{N_a \to\infty}N^{1-r}_a \sum_{j\ge0}j^r\rho_n(j),\quad0\le r\le3.
\end{equation*}
At a non-grid time one may use either adjacent grid time: both approach the same positive time, and their locally uniform limits agree. Equivalently, temporal interpolation gives a convex combination of these discrete measures. Spatial interpolation does not change the limit, by the compact Riemann-sum estimate and the same tail control. The bounds in \eqref{eq:Amoment} follow from \eqref{eq:weightedmass} and \eqref{eq:lowmoments}; the first-moment identity is
\[
 \sum_jj\rho_n(j)=1-\sum_j\rho_n(j)=1-(m-1)\theta_n\longrightarrow1.
\]

For the more precise tail, \eqref{eq:thirdtail} gives, before taking limits,
\[
 N^{-2}_a \sum_{j>\delta N_a}j^3\rho_n(j)
 \le C\left(\frac{n+1}{N_a}\right)^2
 \left\{\Pp(K^{\sharp}_0 > \delta N_a/2)+\frac n{\delta N_a}\right\}.
\]
Portmanteau with bounded lower-semicontinuous truncations, followed by monotone convergence, proves \eqref{eq:limitthirdtail}. The initial tail disappears only as $N_a \to\infty$. For $t$ in a fixed compact interval, the right side is uniformly $O(1/\delta)$.

Now iterate \eqref{eq:rhor} between $\ell=\floor{N_a t_0}$ and $n=\floor{N_a t}$. The product of linear coefficients tends uniformly to one because $\sum_{s=\ell}^{n-1}(1-c_s)=O(N^{-1}_a)$. The quadratic coefficient tends to $1/2$. On compact space-time sets with positive lower time endpoint,
\[
 N^3_a (\rho_s*\rho_s)(k)
 =\frac{1}{N_a} \sum_{i=0}^k[N^2_a \rho_s(i)][N^2_a \rho_s(k-i)]
\]
converges locally uniformly to the corresponding continuum convolution. For $r\ge3$, \eqref{eq:convnorms} gives a total contribution after time summation and multiplication by $N^2_a$ bounded by $CN^{2-r}_a$; each such contribution tends to zero. Fixed right shifts $T^{r-2}$ do not alter that bound. The quadratic space and time Riemann sums therefore give \eqref{eq:mild}. Establishing it first on a countable dense set and using continuity yields every $t_0,t,x$ as stated.
\end{proof}

The pointwise bound also passes to the limit as $u(t,x)\le Ct(t+x)^{-3}$. This bound alone would not control the third-moment tail, so Lemma~\ref{lem:spine} is needed to pass the third moments to the limit.

\subsection{Proof of Lemma~\ref{lem:momentODE}} \label{subsec:prop:momentODE_proof}
Equation \eqref{eq:mild} gives the time regularity and moment identities needed in Section \ref{sec:rigidity}. We derive the moment equations in weak form, using the continuous boundary trace and the tail estimate from Proposition~\ref{prop:limit}.

\begin{lemma}\label{lem:momentODE}
The functions $A_0,A_2,A_3$ and, for each $p\ge0$,
\[
 U(t,p)=\int_0^\infty e^{-px}u(t,x)\dd x
\]
are locally absolutely continuous. With $\beta(t)=u(t,0)$,
\begin{align}
 & A_0'=-\beta+\tfrac12 A_0^2, \quad
 A_2'=A_0A_2-1, \quad A_3'=A_0A_3,\label{eq:momentODE}\\
 & \partial_tU(t,p)=pU(t,p)+\tfrac12U(t,p)^2-\beta(t).\label{eq:laplaceODE}
\end{align}
\end{lemma}
\begin{proof}
The mild equation implies $u_t=u_x+\tfrac12u*u$ in distributions. More explicitly, for $\varphi\in C_c^1([0,\infty))$, translating in \eqref{eq:mild} and integrating gives the following identity in the sense of distributions in time:
\begin{multline} \label{eq:weak_identity}
 \frac{\mathrm d}{\mathrm dt}\int\varphi(x)u(t,x)\dd x
 = I_1 + I_2, \quad \text{where} \\
 I_1 = -\varphi(0)u(t,0)-\int\varphi'(x)u(t,x)\dd x, \quad I_2 =
 +\frac12\iint\varphi(x+y)u(t,x)u(t,y)\dd x\dd y.
\end{multline}
Continuity of $u$ up to zero identifies the boundary term. All terms are locally integrable in time by \eqref{eq:Amoment}.

To use $1,x^2,x^3$ and $e^{-px}$ as test functions in \eqref{eq:weak_identity}, we approximate them by smooth, compactly supported functions. Let $\psi$ be a smooth function on $[0,\infty)$ such that $\psi(x)=1$ for $x\le1$, $\psi(x)=0$ for $x>2$, and $\psi$ is decreasing on $[1,2]$. For $R>0$, put $\psi_R(x)=\psi(x/R)$ and replace each test function $f$ by $f\psi_R$. We prove \eqref{eq:momentODE} below; the proof of \eqref{eq:laplaceODE} is similar. For $f(x)=x^r$ with $r=1,2,3$, the derivatives of $f\psi_R$ are bounded by $C_r x^{r-1}$, uniformly in $R$. For $f(x)=1$, they are bounded by $C/R$ and supported in $[R,2R]$. In all cases, the term $I_1$ is controlled by moments of order at most two, while the integrand in $I_2$ is bounded by $C_r(x^r+y^r)$ and is integrable by \eqref{eq:Amoment}. The left-hand sides converge uniformly on every compact time interval contained in $(0,\infty)$ by \eqref{eq:limitthirdtail}. Taking $R\to\infty$ proves absolute continuity and the stated differential identities. In the convolution moments, $A_1=1$ cancels the $A_2$ term in the third equation and gives $A_2'=A_0A_2-1$. Local continuity of the moments follows from continuity of $u$ and the uniform tail bound on such intervals. The right-hand sides of these equations are therefore continuous, so the absolutely continuous functions are in fact continuously differentiable. No differentiability of $u$ in space is required.
\end{proof}

\section{Proof of Proposition~\ref{prop:identification}}\label{sec:rigidity}
Recall the functions $A_1$, $A_2$, $A_3$ and $U(t,p)$ defined in Sections~\ref{subsec:prop:limit_proof} and~\ref{subsec:prop:momentODE_proof}. Fix a subsequential limit $u$. Recall that its first moment is one, its third moment is $O(t^2)$, and its third-moment tail satisfies \eqref{eq:limitthirdtail}. Using these estimates, we show that \eqref{eq:momentODE} and \eqref{eq:laplaceODE} uniquely determine $u(t,x)$ as the function in \eqref{eq:identifiedprofile}. Our technical tools are from \cite{Kotani}.

\subsection{Construction of the string of entrance type}
Fix a normalization constant $\kappa>0$ and put
\begin{equation*}
 Y(t)=\frac{\kappa}{A_3(t)}.
\end{equation*}
Then $Y'=-A_0Y$, $Y(t)\asymp t^{-2}$, and $(A_2Y)'=-Y$. Since $A_2(t)Y(t)=O(t^{-1})\to0$ as $t\to\infty$,
\[
 A_2(t)Y(t)=\int_t^\infty Y(s)\dd s.
\]
Define the increasing map $\xi:(0,\infty)\to(-\infty,0)$ and the map $\mathfrak m:(-\infty,0)\to(0,\infty)$ by
\begin{equation*}
 \xi(t)=-\int_t^\infty Y(s)\dd s=-A_2(t)Y(t),\qquad
 \mathfrak m(\xi(t))=\int_0^tY(s)^{-1}\dd s.
\end{equation*}
We have
\begin{equation*}
 \xi'(t)=Y(t),\qquad \mathfrak m'(\xi(t))=Y(t)^{-2}.
\end{equation*}
Extend $\mathfrak m(\xi)=\infty$ for $\xi\ge0$. The estimates above give
\[
 -\xi(t)\asymp t^{-1},\qquad \mathfrak m(\xi(t))\asymp t^3,
 \qquad \mathfrak m'(\xi)\asymp|\xi|^{-4}\quad(\xi\to-\infty).
\]
Consequently $\mathfrak m(-\infty)=0$, the support is nonempty, the finite right endpoint is zero, and for every $b<0$,
\begin{equation}\label{eq:stringintegrability}
 \int_{-\infty}^b|\xi|\dd\mathfrak m(\xi)<\infty,
 \qquad \int_{-\infty}^b\xi^2\dd\mathfrak m(\xi)<\infty.
\end{equation}
The first condition in \eqref{eq:stringintegrability} is Kotani's entrance condition \cite[Condition~(2.1)]{Kotani}; thus $\mathfrak m$ is a \emph{string of entrance type} in his terminology. The second is the stronger condition used below to define the characteristic function $h(\lambda)$ associated with the string. Also $\mathfrak m(0-)=\int_0^\infty Y^{-1}\dd t=\infty$, so the spectral measure $\sigma$ of the string $\mathfrak{m}$ is unique by \cite[equation (2.3)]{Kotani}.

In the following, we use the arguments and results of \cite{Kotani}. For nondecreasing strings of entrance type satisfying the first condition in \eqref{eq:stringintegrability}, equality of spectral measures implies equality of the strings up to translation \cite[Theorem 1]{Kotani}. Under the stronger second condition, we follow \cite[Section~1]{Kotani} and define the characteristic function $h$ associated with the singular string $\mathfrak m$ by
\begin{equation}\label{eq:characteristic}
 h(\lambda)=\lim_{\xi\to-\infty}\left\{\xi+\varphi_\lambda(\xi)
 \int_\xi^0\varphi_\lambda(v)^{-2}\dd v\right\} = a + \int^{\infty}_0 \left( \frac{1}{\xi - \lambda} - \frac{\xi}{\xi^2 + 1} \right) \dd \sigma(\xi),\qquad\lambda<0,
\end{equation}
where $\sigma$ is the spectral measure and $\varphi_\lambda$ is the unique solution of 
\begin{equation} \label{eq:entrance_solution}
    \varphi(\xi) = 1 - \lambda \int^{\xi}_{-\infty} (\xi - v) \varphi(v) \dd \mathfrak{m}(v),
\end{equation}
with $\varphi_\lambda(-\infty)=1$. Therefore equality of these characteristic functions for every negative parameter implies equality of the spectral measures. Next, fixing the right endpoint of $\mathfrak{m}(\xi)$ at zero removes the translation freedom. We finally determine $u(t, x)$ uniquely from $\mathfrak{m}(\xi)$.

\subsection{Notation and estimates}
By H\"older and Cauchy--Schwarz,
\begin{equation}\label{eq:A3lower}
 1=A_1^3\le A_0^2A_3,\qquad A_2^2\le A_1A_3=A_3.
\end{equation}
Thus $ct^2\le A_3(t)\le Ct^2$ and $A_2(t)\le Ct$. For $p>0$, set
\[
 Z(t,p)=p-A_0(t)+U(t,p)=\int(e^{-px}-1+px)u(t,x)\dd x.
\]
Then $Z>0$, $Z\le p^2A_2/2=O_p(t)$, and cancellation of the boundary term in \eqref{eq:laplaceODE} gives
\begin{equation}\label{eq:ZODE}
 \partial_tZ=A_0Z+\tfrac12(Z^2-p^2).
\end{equation}

\begin{lemma}\label{lem:remainder}
For each fixed $p>0$,
\begin{equation*}
 Z(t,p)=\frac{p^2}{2}A_2(t)-\frac{p^3}{6}A_3(t)+o(A_3(t))\qquad(t\downarrow0).
\end{equation*}
In fact, the remainder divided by $A_3(t)$ is $O_p(\sqrt t)$ for $0<t\le1$.
\end{lemma}
\begin{proof}
The nonnegative remainder
$R_p(x)=e^{-px}-1+px-p^2x^2/2+p^3x^3/6$ satisfies
\[
 0\le R_p(x)\le \min\{p^4x^4/24,\ p^3x^3/6\}.
\]
Split its integral at a number $\delta>0$. By \eqref{eq:limitthirdtail} and \eqref{eq:A3lower},
\[
 \frac{\int R_p(x)u(t,x)\dd x}{A_3(t)}
 \le C_p\delta+C_p\frac{t}{\delta}.
\]
Taking $\delta=\sqrt t$ proves the assertion. The quantitative tail in \eqref{eq:limitthirdtail} is the input that makes the normalized remainder vanish.
\end{proof}


For $p>0$ put
\begin{equation*}
 \Phi(t,p)=\exp\left\{-\frac12\int_0^tZ(s,p)\dd s\right\},\qquad
 f_p(\xi(t))=\frac2{p^2}Y(t)Z(t,p)\Phi(t,p).
\end{equation*}
The bounds on $Z$ imply $0<\Phi\le1$ and $\Phi=1+O_p(t^2)$ near zero. Direct differentiation using \eqref{eq:ZODE} gives
\begin{equation}\label{eq:fODE}
 \frac{\mathrm d}{\mathrm dt}f_p(\xi(t))=-Y(t)\Phi(t,p),\qquad
 f'_p(\xi(t))=-\Phi(t,p),\qquad
 f''_p=\frac{p^2}{4}\mathfrak m'f_p.
\end{equation}
The solution is positive, decreasing, and satisfies $f_p(0-)=0$, since
$f_p(\xi(t))\le Y(t)A_2(t)=O(t^{-1})$ as $t\to\infty$.

Lemma~\ref{lem:remainder}, together with $YA_3=\kappa$, gives
\begin{align}\label{eq:fpentrance}
 f_p(\xi(t))+\xi(t)
 &=YA_2(\Phi-1)-\frac{\kappa p}{3}\Phi
   +\frac2{p^2}Y\Phi\int R_p(x)u(t,x)\dd x\notag\\
 &=-\frac{\kappa p}{3}+O_p(\sqrt t).
\end{align}
Furthermore, $f'_p(\xi)\to-1$ as $\xi\to-\infty$.

\subsection{Computation of \texorpdfstring{$\mathfrak{m}$}{m} and \texorpdfstring{$u(t, x)$}{u(t, x)}}
Let $\lambda = -p^2/4$ and let $\varphi_{\lambda}$ be defined as the solution of \eqref{eq:entrance_solution}. Thus $\varphi_{\lambda} \ge1$ and $\varphi'_{\lambda} \ge0$. Since
$B(\xi)=\int_{-\infty}^\xi(\xi-v)\mathfrak m'(v)\dd v=O(|\xi|^{-2})$, monotonicity gives
$\varphi_{\lambda}(\xi)\le1 - \lambda B(\xi)\varphi_{\lambda}(\xi)$. For sufficiently large negative $\xi$,
\[
 \varphi_{\lambda} (\xi)=1+O_{\lambda} (|\xi|^{-2}),\qquad
 \varphi'_{\lambda} (\xi)=O_{\lambda} (|\xi|^{-3}).
\]
Equations \eqref{eq:fODE} and \eqref{eq:fpentrance} then imply that the constant Wronskian is
\[
 f_p\varphi_{\lambda} '-f'_p\varphi_{\lambda} = 1.
\]
Since $f_p/\varphi_{\lambda} \to0$ at $0-$, integration of $(f_p/\varphi_{\lambda})'=-\varphi^{-2}_{\lambda}$ proves
\[
 f_p(\xi)=\varphi_{\lambda} (\xi)\int_\xi^0\varphi(v)^{-2}_{\lambda} \dd v.
\]
Substituting this identity into \eqref{eq:characteristic} expresses $h(\lambda)$ as the limit of $f_p(\xi)+\xi$ as $\xi\to-\infty$. Hence \eqref{eq:fpentrance} gives
\begin{equation*}
 h(\lambda)=-\frac{\kappa p}{3} = -\frac{2\kappa}{3} \sqrt{-\lambda}.
\end{equation*}
This is an identity on the whole negative axis.

On the other hand, put $d=2\kappa/3$ and consider the explicit string
\begin{equation*}
 \mathfrak m_*(\xi)=\frac{d^2}{3(-\xi)^3}\quad(\xi<0),\qquad
 \mathfrak m_*(\xi)=\infty\quad(\xi\ge0).
\end{equation*}
A direct computation gives the solution of \eqref{eq:entrance_solution} with $\mathfrak m$ replaced by $\mathfrak m_*$ and the corresponding characteristic function: for every $\lambda<0$,
\begin{equation*}
    \varphi_{*, \lambda} = \frac{\sinh(d\sqrt{-\lambda}/(-\xi))}{d\sqrt{-\lambda}/(-\xi)}, \qquad h_*(\lambda) = -d\sqrt{-\lambda}.
\end{equation*}
Since $h_*(\lambda) = h(\lambda)$, we have that $\mathfrak{m}_*$ and $\mathfrak{m}$ differ only by a translation. Kotani's uniqueness theorem and their common finite right endpoint give $\mathfrak m=\mathfrak m_*$. Moreover, we also find that
\begin{equation*}
    f_p(\xi) = (-\xi) \exp \left( \frac{dp}{2\xi} \right).
\end{equation*}

To compute $u(t,x)$, we compare the densities of $\mathfrak m$ and $\mathfrak m_*$:
\[
 Y(t)^{-2}=d^2\xi(t)^{-4},\qquad \xi'(t)=Y(t).
\]
Thus $\xi'=\xi^2/d$. The condition $\xi(t)\to-\infty$ as $t\downarrow0$ fixes the integration constant, giving
\[
 \xi(t)=-d/t,\qquad Y(t)=d/t^2,\qquad A_0(t)=-Y'(t)/Y(t)=2/t.
\]
From $f'_p(\xi(t))=-\Phi(t,p)$ we obtain
\[
 \Phi(t,p)=(1+pt/2)e^{-pt/2},\qquad
 Z(t,p)=-2\partial_t\log\Phi(t,p)=\frac{tp^2}{tp+2}.
\]
Therefore
\begin{equation*}
 U(t,p)=A_0(t)-p+Z(t,p)=\frac4{t(tp+2)},\qquad
 u(t,x)=\frac4{t^2}e^{-2x/t}.
\end{equation*}
Uniqueness of Laplace transforms proves the last equality. The time variable is the original scaling $t=n/N$. This proves Proposition~\ref{prop:identification}.

\appendix
\section{Sharpness of the moment assumption for the survival asymptotic}\label{sec:sharpness}
In this appendix, we show that, for any fixed $r\in[0,3)$, the assumption $\E[X_0^rm^{X_0}]<\infty$ does not, in general, guarantee the survival asymptotic $n^2p_n\to4/(m-1)^2$.

\begin{theorem}\label{thm:sharpness}
For every fixed $m\ge2$ and every $r\in[0,3)$, there exists a nondegenerate critical initial law with $\E[X_0^rm^{X_0}]<\infty$ but $\E[X_0^3m^{X_0}]=\infty$ for which $n^2p_n$ does not converge to $4/(m-1)^2$.
\end{theorem}

The proof constructs a critical law for which $\Pp(X_0=k)$ is proportional to $m^{-k}k^{-\alpha}$ for $k\ge2$, where $\alpha\in(\max\{3,r+1\},4)$. The choice $\alpha>r+1$ ensures that $\E[X_0^rm^{X_0}]<\infty$. The product estimate of Chen and Shi~\cite[Theorem~1.1]{CS} gives
\[
 \prod_{j=0}^{n-1}G_j^{m-1}\asymp n^{\alpha-2}.
\]
The following lemma shows that, for a critical initial law, the survival asymptotic $n^2p_n\to4/(m-1)^2$ implies $\prod_{j=0}^{n-1}G_j^{m-1}=n^{2+o(1)}$, without assuming $\E[X_0^3m^{X_0}]<\infty$. This contradicts the preceding estimate since $\alpha-2<2$.

\begin{lemma}\label{lem:productasymptotic}
Fix $m\ge2$ and a nonnegative integer-valued initial law satisfying
\[
 (m-1)\E[X_0m^{X_0}]=\E[m^{X_0}]<\infty.
\]
Let $(X_n)$ evolve according to \eqref{eq:model}, with $p_n$ and $G_n$
defined by \eqref{eq:notation}.
Suppose that
\begin{equation}\label{eq:survivalassumption}
 n^2p_n\longrightarrow\frac4{(m-1)^2}.
\end{equation}
Then
\begin{equation}\label{eq:productasymptotic}
 \prod_{j=0}^{n-1}G_j^{m-1}=n^{2+o(1)}.
\end{equation}
\end{lemma}
\begin{proof}
We first show that
\begin{equation}\label{eq:productlimitassumption}
 n(G_n-1)\longrightarrow\frac2{m-1},
\end{equation}
and then use this limit to prove \eqref{eq:productasymptotic}.

\textbf{Proof of \eqref{eq:productlimitassumption}.}
Put $\varepsilon_n=G_n-1$. The recursion \eqref{eq:pgf} preserves
criticality, so the tilted law $q_n$ has mean $1/(m-1)$.
Since $G_n\ge1$, for every integer $K\ge1$,
\[
 \theta_n=\sum_{k\ge1}q_{n,k}
 \le m^Kp_n+\frac1{K+1}\sum_{k>K}kq_{n,k}
 \le m^Kp_n+\frac1{(m-1)(K+1)}.
\]
By \eqref{eq:survivalassumption}, $p_n\to0$. Letting first $n\to\infty$
and then $K\to\infty$ gives $\theta_n\to0$. Hence
\[
 G_n=\frac{1-p_n}{1-\theta_n}\longrightarrow1,
\]
so $\varepsilon_n\to0$.

Evaluating \eqref{eq:pgf} at $m$ gives
\[
 \varepsilon_{n+1}=F_n(\varepsilon_n),\qquad
 F_n(x)=\frac{(1+x)^m+(m-1)(1-p_n)^m}{m}-1.
\]
We now compare $\varepsilon_n$ with $b/n$. For each fixed $b>0$,
\eqref{eq:survivalassumption} and a Taylor expansion give
\begin{equation}\label{eq:survivalbarrier}
 F_n(b/n)-\frac b{n+1}
 =\frac1{n^2}\left(\frac{m-1}{2}b^2+b-\frac4{m-1}+o(1)\right).
\end{equation}
The polynomial in parentheses has the unique positive root
$b_*=2/(m-1)$, and $F_n'(x)=(1+x)^{m-1}\ge1$ for $x\ge0$.
If $b>b_*$, the left-hand side of \eqref{eq:survivalbarrier} is positive
for all $n\ge N_b$, where $N_b\ge1$. Suppose that
$\varepsilon_{n_0}\ge b/n_0$ for an index $n_0\ge N_b$.
Whenever $n\ge n_0$ and $\varepsilon_n\ge b/n$, the derivative bound gives
\[
\begin{aligned}
 \varepsilon_{n+1}-\frac b{n+1}
 &=F_n(\varepsilon_n)-F_n(b/n)
   +F_n(b/n)-\frac b{n+1}\\
 &>\varepsilon_n-\frac bn\ge0.
\end{aligned}
\]
By induction, the difference is strictly positive from $n_0+1$ onward
and remains nondecreasing. This contradicts
$\varepsilon_n-b/n\to0$. Thus $\varepsilon_n<b/n$ eventually.

If $0<b<b_*$, choose $N_b$ so that the left-hand side of
\eqref{eq:survivalbarrier} is negative for all $n\ge N_b$.
Suppose that $\varepsilon_{n_0}\le b/n_0$ for some $n_0\ge N_b$.
Whenever $n\ge n_0$ and $0\le\varepsilon_n\le b/n$, we have
\[
\begin{aligned}
 \frac b{n+1}-\varepsilon_{n+1}
 &=F_n(b/n)-F_n(\varepsilon_n)
   +\frac b{n+1}-F_n(b/n)\\
 &>\frac bn-\varepsilon_n\ge0.
\end{aligned}
\]
Induction again gives a difference that is strictly positive from
$n_0+1$ onward and nondecreasing. But this difference is at most $b/n$
because $\varepsilon_n\ge0$, a contradiction as $n\to\infty$.
Thus $\varepsilon_n>b/n$ eventually.
Letting $b\downarrow b_*$ and $b\uparrow b_*$ in these two bounds proves
\eqref{eq:productlimitassumption}.

\textbf{Proof of \eqref{eq:productasymptotic}.}
By \eqref{eq:productlimitassumption}, $G_j-1=O(j^{-1})$, so
$\sum_{j\ge0}(G_j-1)^2<\infty$. Expanding $\log G_j$ at one, we obtain
\[
 \log\left(\prod_{j=0}^{n-1}G_j^{m-1}\right)
 =(m-1)\sum_{j=0}^{n-1}(G_j-1)+O(1)
 =2\log n+o(\log n).
\]
Exponentiating proves the claim.
\end{proof}

\begin{proof}[Proof of Theorem~\ref{thm:sharpness}]
Fix $m\ge2$ and $r\in[0,3)$. Choose $\alpha\in(\max\{3,r+1\},4)$ and $c>0$ sufficiently small. Define a tilted probability law by
\[
 q_{0,k}=ck^{-\alpha}\ (k\ge2),\quad
 q_{0,1}=\frac1{m-1}-c\sum_{k\ge2}k^{1-\alpha},\quad
 q_{0,0}=1-\frac1{m-1}+c\sum_{k\ge2}(k-1)k^{-\alpha}.
\]
All masses are nonnegative, $q_{0,0}>0$, the total mass is one, and the mean is $1/(m-1)$. Put
\[
 G_0=\left(\sum_{k\ge0}m^{-k}q_{0,k}\right)^{-1},\qquad
 \Pp(X_0=k)=G_0m^{-k}q_{0,k}.
\]
The original law is exactly critical, nondegenerate, and has positive mass above one. For $s>0$, $\E[X_0^sm^{X_0}]<\infty$ if and only if $s<\alpha-1$; moreover, $\E[m^{X_0}]<\infty$. Hence $\E[X_0^rm^{X_0}]<\infty$ and $\E[X_0^3m^{X_0}]=\infty$.

The product estimate of Chen and Shi~\cite[Theorem~1.1]{CS} applies because $\Pp(X_0=k)=G_0c\,m^{-k}k^{-\alpha}$ for $k\ge2$. It gives
\[
 \prod_{j=0}^{n-1}G_j^{m-1}\asymp n^{\alpha-2}.
\]
Since $\alpha-2<2$, this is incompatible with the conclusion of
Lemma~\ref{lem:productasymptotic}. Hence $n^2p_n$ cannot converge to
$4/(m-1)^2$, proving the theorem.
\end{proof}
\section*{Funding}
YZ is supported by the Fundamental Research Funds for the Central Universities.

% Candidate funding acknowledgment, pending confirmation that these grants
% apply to the present work. Source: the Acknowledgements in
% https://arxiv.org/html/2307.03720v5 (23 October 2025).
% The supplied four-author manuscript contains no funding statement.
% \section*{Funding}
% Dong Wang acknowledges partial support from NSFC grants 12271502 and
% 11871425 and UCAS start-up grant 118900M043.

\begin{thebibliography}{99}
\bibitem{CDHLS} X. Chen, B. Derrida, Y. Hu, M. Lifshits and Z. Shi (2019).
A max-type recursive model: some properties and open questions.
In \emph{Sojourns in Probability Theory and Statistical Physics--III},
Springer Proceedings in Mathematics \& Statistics \textbf{300}, 166--186.
\href{https://arxiv.org/abs/1705.04787v2}{arXiv:1705.04787v2}.

\bibitem{CHS} X. Chen, Y. Hu and Z. Shi (2022).
The sustainability probability for the critical Derrida--Retaux model.
\emph{Probab. Theory Related Fields} \textbf{182}, 641--684.
\href{https://doi.org/10.1007/s00440-021-01091-z}{doi:10.1007/s00440-021-01091-z}.

\bibitem{CS} X. Chen and Z. Shi (2021).
The stable Derrida--Retaux system at criticality.
In \emph{In and Out of Equilibrium 3: Celebrating Vladas Sidoravicius},
Progress in Probability \textbf{77}, 239--264.
\href{https://arxiv.org/abs/2006.06140v1}{arXiv:2006.06140v1}.

\bibitem{CEGM} P. Collet, J.-P. Eckmann, V. Glaser and A. Martin (1984).
Study of the iterations of a mapping associated to a spin glass model.
\emph{Comm. Math. Phys.} \textbf{94}, 353--370.
\href{https://doi.org/10.1007/BF01224830}{doi:10.1007/BF01224830}.

\bibitem{DR} B. Derrida and M. Retaux (2014).
The depinning transition in presence of disorder: a toy model.
\emph{J. Stat. Phys.} \textbf{156}, 268--290.
\href{https://arxiv.org/abs/1401.6919}{arXiv:1401.6919}.

\bibitem{Kotani} S. Kotani (2013).
Krein's strings whose spectral functions are of polynomial growth.
\emph{Kyoto J. Math.} \textbf{53}, 787--814.
\href{https://arxiv.org/abs/1304.6786v1}{arXiv:1304.6786v1}.

\bibitem{CDDHLS} X. Chen, V. Dagard, B. Derrida, Y. Hu, M. Lifshits and Z. Shi (2021).
The Derrida--Retaux conjecture on recursive models.
\emph{Ann. Probab.} \textbf{49}, 637--670.
\href{https://doi.org/10.1214/20-AOP1457}{doi:10.1214/20-AOP1457}.
\href{https://arxiv.org/abs/1907.01601v3}{arXiv:1907.01601v3}.

\bibitem{CDDS} X. Chen, V. Dagard, B. Derrida and Z. Shi (2020).
The critical behaviors and the scaling functions of a coalescence equation.
\emph{J. Phys. A: Math. Theor.} \textbf{53}, 195202.
\href{https://arxiv.org/abs/2001.00853}{arXiv:2001.00853}.

\bibitem{HMP} Y. Hu, B. Mallein and M. Pain (2020).
An exactly solvable continuous-time Derrida--Retaux model.
\emph{Comm. Math. Phys.} \textbf{375}, 605--651.
\href{https://doi.org/10.1007/s00220-019-03465-w}{doi:10.1007/s00220-019-03465-w}.

\bibitem{LZ} Z. Li and R. Zhang (2025).
Asymptotic behavior of the generalized Derrida--Retaux recursive model.
\emph{J. Stat. Phys.} \textbf{192}, article 98.
\href{https://doi.org/10.1007/s10955-025-03480-3}{doi:10.1007/s10955-025-03480-3}.

\bibitem{AHM} G. Alsmeyer, Y. Hu and B. Mallein (2025).
The Derrida--Retaux model on a geometric Galton--Watson tree.
Preprint.
\href{https://arxiv.org/abs/2502.02991}{arXiv:2502.02991}.

\bibitem{DS} B. Derrida and Z. Shi (2020).
Results and conjectures on a toy model of depinning.
\emph{Moscow Math. J.} \textbf{20}, 695--709.
\href{https://arxiv.org/abs/2005.10208}{arXiv:2005.10208}.

\bibitem{HS} Y. Hu and Z. Shi (2018).
The free energy in the Derrida--Retaux recursive model.
\emph{J. Stat. Phys.} \textbf{172}, 718--741.
\href{https://doi.org/10.1007/s10955-018-2066-1}{doi:10.1007/s10955-018-2066-1}.

\bibitem{CHSdual} X. Chen, Y. Hu and Z. Shi (2024).
The dual Derrida--Retaux conjecture.
\emph{Stochastic Process. Appl.} \textbf{171}, article 104332.
\href{https://arxiv.org/abs/2306.12717}{arXiv:2306.12717}.

\bibitem{ChenSmooth} X. Chen (2026).
Infinite differentiability of the free energy for a Derrida--Retaux system.
\emph{Acta Math. Sin. (Engl. Ser.)} \textbf{42}, 1457--1480.
\href{https://doi.org/10.1007/s10114-026-4249-z}{doi:10.1007/s10114-026-4249-z}.

\bibitem{LZscale} Z. Li and R. Zhang (2024).
Derrida--Retaux type models and related scaling limit theorems.
Preprint.
\href{https://arxiv.org/abs/2411.12189v2}{arXiv:2411.12189v2}.

\bibitem{DDS} B. Derrida, T. Duquesne and Z. Shi (2024).
The critical tree of a renormalization model as a growth-fragmentation process.
\emph{Ann. Inst. H. Poincar\'e Probab. Statist.} \textbf{60}, 1985--2024.
\href{https://doi.org/10.1214/23-AIHP1390}{doi:10.1214/23-AIHP1390}.

\bibitem{DuS} T. Duquesne and Z. Shi (2026).
The continuous Derrida--Retaux branching process in the Brownian CRT.
Preprint.
\href{https://arxiv.org/abs/2609.11435v1}{arXiv:2609.11435v1}.

\bibitem{Yaglom} A. M. Yaglom (1947).
Certain limit theorems of the theory of branching random processes.
\emph{Doklady Akad. Nauk SSSR (N.S.)} \textbf{56}, 795--798.
MR22045.

\end{thebibliography}
\end{document}
